\section{Forma de Lur'e, transferencia y función descriptiva}
\label{sec:lure_transferencia}

\HAFOCardLure

\subsection{Separación común de los dos sistemas}

Los sistemas de \cref{eq:chua_no_suave,eq:chua_arctan} tienen la misma
estructura. Se escribe
\[
X=(x,y,z)^{\T},\qquad
\sigma=r^{\T}X=x,\qquad
h(x)=\ell x+\psi(x),
\]
y se separa la parte lineal de la no linealidad escalar:
\begin{equation}
\begin{aligned}
{}^{\mathrm C}D_{t_0}^{q}X
&=PX+b\,\psi(\sigma),&
\sigma&=r^{\T}X,\\
P&=
\begin{pmatrix}
-\alpha(1+\ell)&\alpha&0\\
1&-1&1\\
0&-\beta&-\gamma
\end{pmatrix},&
b&=
\begin{pmatrix}
-\alpha\\0\\0
\end{pmatrix},&
r&=
\begin{pmatrix}
1\\0\\0
\end{pmatrix}.
\end{aligned}
\label{eq:lure_comun}
\end{equation}
La correspondencia con cada modelo es
\begin{equation}
\begin{array}{c|c|c}
\text{modelo}&\ell&\psi(\sigma)\\ \hline
\text{Chua no suave}
&m_1&(m_0-m_1)\operatorname{sat}(\sigma)\\[1mm]
\text{Chua con arctangente}
&a_1&a_2\arctan(\rho\sigma).
\end{array}
\label{eq:correspondencia_lure}
\end{equation}
La equivalencia se obtiene por comparación componente a componente:
\begingroup
\def\HAFOderivationpart{1}
% Este archivo es una biblioteca de derivaciones. Cada bloque se inserta en el
% lugar del cuerpo donde se usa, mediante \HAFOderivationpart. No se imprime
% como apéndice independiente.
\ifnum\HAFOderivationpart=1\relax
\subsection{Derivación completa de la forma de Lur'e}
\label{app:verificacion_lure}

Sea \(h(x)=\ell x+\psi(x)\). La primera ecuación de los dos modelos es
\[
{}^{\mathrm C}D_{t_0}^{q}x
=\alpha(y-x-h(x)).
\]
Al distribuir \(\alpha\),
\begin{align}
\alpha(y-x-h(x))
&=\alpha y-\alpha x-\alpha\ell x-\alpha\psi(x)\nonumber\\
&=-\alpha(1+\ell)x+\alpha y-\alpha\psi(x).
\label{eq:expansion_primera_fila}
\end{align}
Por otra parte,
\[
PX=
\begin{pmatrix}
-\alpha(1+\ell)x+\alpha y\\
x-y+z\\
-\beta y-\gamma z
\end{pmatrix},
\qquad
b\psi(r^{\T}X)=
\begin{pmatrix}
-\alpha\psi(x)\\0\\0
\end{pmatrix}.
\]
La suma reproduce \cref{eq:expansion_primera_fila} y las otras dos
ecuaciones. Para la saturación se toma
\(\ell=m_1\) y
\(\psi=(m_0-m_1)\operatorname{sat}\); para la arctangente,
\(\ell=a_1\) y \(\psi=a_2\arctan(\rho\,\cdot)\).

\fi
\ifnum\HAFOderivationpart=2\relax
\subsection{Derivación completa del determinante y la transferencia}
\label{app:transferencia_explicita}

Con \(\lambda=s^q\),
\begin{equation}
\lambda I-P=
\begin{pmatrix}
\lambda+\alpha(1+\ell)&-\alpha&0\\
-1&\lambda+1&-1\\
0&\beta&\lambda+\gamma
\end{pmatrix}.
\label{eq:matriz_resolvente}
\end{equation}
El menor de la posición \((1,1)\) es
\begin{align}
Q(\lambda)
&=
\det
\begin{pmatrix}
\lambda+1&-1\\
\beta&\lambda+\gamma
\end{pmatrix}\nonumber\\
&=(\lambda+1)(\lambda+\gamma)+\beta.
\label{eq:menor_Q}
\end{align}
Al desarrollar \(\det(\lambda I-P)\) por la primera fila,
\begin{align}
\Delta_\ell(\lambda)
&=\bigl[\lambda+\alpha(1+\ell)\bigr]Q(\lambda)
-(-\alpha)
\det
\begin{pmatrix}
-1&-1\\
0&\lambda+\gamma
\end{pmatrix}\nonumber\\
&=\bigl[\lambda+\alpha(1+\ell)\bigr]Q(\lambda)
-\alpha(\lambda+\gamma).
\label{eq:determinante_transferencia}
\end{align}

Para calcular la primera componente de
\((\lambda I-P)^{-1}b\), la regla de Cramer sustituye la primera columna de
\cref{eq:matriz_resolvente} por
\(b=(-\alpha,0,0)^{\T}\). El determinante del numerador es
\[
\det
\begin{pmatrix}
-\alpha&-\alpha&0\\
0&\lambda+1&-1\\
0&\beta&\lambda+\gamma
\end{pmatrix}
=-\alpha Q(\lambda).
\]
Como \(r^{\T}\) selecciona esa primera componente,
\begin{equation}
W_q(s)
=r^{\T}(s^qI-P)^{-1}b
=-\alpha\frac{Q(s^q)}{\Delta_\ell(s^q)}.
\label{eq:transferencia_cramer}
\end{equation}

La identidad
\[
P-s^qI=-(s^qI-P)
\]
implica
\[
r^{\T}(P-s^qI)^{-1}b=-W_q(s).
\]
Por ello, cambiar el orden de los términos dentro de la resolvente cambia
simultáneamente el signo de la transferencia y el signo del cierre. No cambia
la solución física si ambas modificaciones se realizan de manera coherente.

\fi
\ifnum\HAFOderivationpart=3\relax
\subsection{Derivación completa de los equilibrios}
\label{app:equilibrios}

En un equilibrio \(e=(x_e,y_e,z_e)^{\T}\), las dos últimas ecuaciones de
ambos sistemas dan
\[
0=x_e-y_e+z_e,\qquad
0=-\beta y_e-\gamma z_e.
\]
De la primera,
\[
z_e=y_e-x_e.
\]
Al sustituirla en la segunda,
\begin{align*}
0
&=-\beta y_e-\gamma(y_e-x_e)\\
&=-(\beta+\gamma)y_e+\gamma x_e.
\end{align*}
Si \(\beta+\gamma\neq0\),
\begin{equation}
y_e=\frac{\gamma}{\beta+\gamma}x_e,\qquad
z_e=-\frac{\beta}{\beta+\gamma}x_e.
\label{eq:equilibrios_yz}
\end{equation}
La primera ecuación de equilibrio exige
\[
0=y_e-x_e-h(x_e).
\]
Mediante \cref{eq:equilibrios_yz},
\begin{equation}
h(x_e)+\frac{\beta}{\beta+\gamma}x_e=0.
\label{eq:equilibrio_escalar_comun}
\end{equation}
Esta única ecuación escalar determina todos los equilibrios; las otras dos
coordenadas se recuperan después con \cref{eq:equilibrios_yz}.

\subsubsection{Equilibrios del Chua no suave}

En la región \(|x_e|<1\), \(h(x_e)=m_0x_e\). Entonces
\begin{equation}
\left(m_0+\frac{\beta}{\beta+\gamma}\right)x_e=0.
\label{eq:equilibrio_interior_ns}
\end{equation}
Cuando el coeficiente no es cero, la única raíz interior es \(x_e=0\).

Para \(x_e>1\),
\[
h(x_e)=m_1x_e+(m_0-m_1),
\]
y \cref{eq:equilibrio_escalar_comun} queda
\[
\left(m_1+\frac{\beta}{\beta+\gamma}\right)x_e+(m_0-m_1)=0.
\]
Por tanto,
\begin{equation}
x_\star=
\frac{m_1-m_0}
{m_1+\dfrac{\beta}{\beta+\gamma}}.
\label{eq:raiz_exterior_ns}
\end{equation}
Si \(x_\star>1\), existe el equilibrio positivo. Como la no linealidad es
impar, existe también la raíz \(-x_\star<-1\). Los tres equilibrios son
\begin{equation}
E_0=(0,0,0),\qquad
E_\pm=
\left(
\pm x_\star,\,
\pm\frac{\gamma}{\beta+\gamma}x_\star,\,
\mp\frac{\beta}{\beta+\gamma}x_\star
\right).
\label{eq:equilibrios_ns_simbolicos}
\end{equation}
La condición regional \(x_\star>1\) debe comprobarse; no basta con resolver la
ecuación lineal exterior.

\subsubsection{Equilibrios del Chua con arctangente}

Al usar
\(h(x)=a_1x+a_2\arctan(\rho x)\) en
\cref{eq:equilibrio_escalar_comun}, se obtiene
\begin{equation}
\left(a_1+\frac{\beta}{\beta+\gamma}\right)x_e
+a_2\arctan(\rho x_e)=0.
\label{eq:equilibrio_escalar_arctan}
\end{equation}
La expresión es impar, de modo que \(x_e=0\) siempre es raíz y toda raíz
positiva no nula tiene una compañera negativa. Las raíces no nulas se
determinan numéricamente y se sustituyen en \cref{eq:equilibrios_yz}. Esta
ecuación se evalúa con \(\rho=1\) para el caso bibliográfico y con el valor de
\(\rho\) de \cref{eq:parametros_c590} para el caso \texttt{c590}; no deben
intercambiarse las dos parametrizaciones.

\fi
\ifnum\HAFOderivationpart=4\relax
\subsection{Derivación completa de los jacobianos regionales}
\label{app:jacobianos}

Para la forma de Lur'e,
\[
F(X)=PX+b\psi(r^{\T}X).
\]
La regla de la cadena produce
\begin{equation}
J_e=P+b\,\psi'(r^{\T}e)r^{\T}.
\label{eq:jacobiano_lure}
\end{equation}
Como \(r^{\T}e=x_e\), solamente cambia la entrada \((1,1)\).

En el sistema no suave,
\[
\psi'_{\mathrm{ns}}(x)=
\begin{cases}
m_0-m_1,&|x|<1,\\
0,&|x|>1.
\end{cases}
\]
Por ello,
\begin{equation}
J_{\mathrm{in}}=
\begin{pmatrix}
-\alpha(1+m_0)&\alpha&0\\
1&-1&1\\
0&-\beta&-\gamma
\end{pmatrix},
\qquad
J_{\mathrm{out}}=
\begin{pmatrix}
-\alpha(1+m_1)&\alpha&0\\
1&-1&1\\
0&-\beta&-\gamma
\end{pmatrix}.
\label{eq:jacobianos_ns}
\end{equation}
En \(x=\pm1\) la derivada clásica no existe. Los equilibrios usados en este
reporte no se encuentran sobre esas superficies, por lo que se emplea el
jacobiano de la región que contiene a cada uno.

Para la arctangente,
\[
\psi'_{\mathrm a}(x)
=\frac{a_2\rho}{1+\rho^2x^2},
\]
y
\begin{equation}
J_e=
\begin{pmatrix}
\displaystyle
-\alpha\left(1+a_1+
\frac{a_2\rho}{1+\rho^2x_e^2}\right)
&\alpha&0\\[3mm]
1&-1&1\\
0&-\beta&-\gamma
\end{pmatrix}.
\label{eq:jacobiano_arctan}
\end{equation}
Los valores propios de
\cref{eq:jacobianos_ns,eq:jacobiano_arctan} se clasifican con
\cref{eq:matignon}; la estabilidad local no sustituye el sondeo de cuenca.

\fi
\ifnum\HAFOderivationpart=5\relax
\subsection{Derivación completa del polinomio característico}
\label{app:polinomio_caracteristico}

Sea \(g=h'(x_e)\) la pendiente total de la no linealidad en un equilibrio.
El jacobiano de cualquiera de las dos familias puede escribirse como
\begin{equation}
J(g)=
\begin{pmatrix}
-\alpha(1+g)&\alpha&0\\
1&-1&1\\
0&-\beta&-\gamma
\end{pmatrix}.
\label{eq:jacobiano_pendiente_total}
\end{equation}
Por tanto,
\[
\lambda I-J(g)=
\begin{pmatrix}
\lambda+\alpha(1+g)&-\alpha&0\\
-1&\lambda+1&-1\\
0&\beta&\lambda+\gamma
\end{pmatrix}.
\]
El desarrollo por la primera fila, igual al de
\cref{eq:determinante_transferencia}, produce
\begin{equation}
\chi_g(\lambda)
=\bigl[\lambda+\alpha(1+g)\bigr]
\bigl[(\lambda+1)(\lambda+\gamma)+\beta\bigr]
-\alpha(\lambda+\gamma).
\label{eq:polinomio_caracteristico_agrupado}
\end{equation}
Como
\[
(\lambda+1)(\lambda+\gamma)+\beta
=\lambda^2+(1+\gamma)\lambda+(\beta+\gamma),
\]
la multiplicación y agrupación por potencias de \(\lambda\) da
\begin{equation}
\begin{aligned}
\chi_g(\lambda)
={}&\lambda^3+
\bigl[1+\gamma+\alpha(1+g)\bigr]\lambda^2\\
&+\bigl[
\beta+\gamma+\alpha(1+g)(1+\gamma)-\alpha
\bigr]\lambda\\
&+\alpha\bigl[(1+g)(\beta+\gamma)-\gamma\bigr].
\end{aligned}
\label{eq:polinomio_caracteristico_expandido}
\end{equation}
En el sistema no suave se sustituye \(g=m_0\) en \(E_0\) y \(g=m_1\)
en \(E_\pm\). En el sistema con arctangente,
\[
g=a_1+\frac{a_2\rho}{1+\rho^2x_e^2}.
\]
Las raíces de \cref{eq:polinomio_caracteristico_expandido} son los valores
propios utilizados en \cref{eq:matignon}; así, la clasificación local se
obtiene a partir del mismo campo que define la integración.

\fi
\ifnum\HAFOderivationpart=6\relax
\subsection{Derivación completa para la saturación centrada}
\label{app:df_saturacion_centrada}

Sea
\[
\psi_{\mathrm{ns}}(\sigma)=\Delta\operatorname{sat}(\sigma),
\qquad \Delta=m_0-m_1.
\]
Si \(0<A\leq1\), no hay saturación y
\[
\psi_{\mathrm{ns}}(A\sin\theta)=\Delta A\sin\theta.
\]
Como \(\int_0^\pi\sin^2\theta\dd\theta=\pi/2\),
\[
N_{\mathrm{ns}}(A)
=\frac{2}{\pi A}\Delta A\frac{\pi}{2}
=\Delta.
\]

Para \(A>1\), defínase
\[
\theta_0=\arcsin\left(\frac1A\right),
\qquad 0<\theta_0<\frac{\pi}{2}.
\]
En \(0\leq\theta\leq\theta_0\) y
\(\pi-\theta_0\leq\theta\leq\pi\), la entrada permanece en la parte lineal;
en el intervalo central la saturación vale \(1\). Entonces
\begin{align}
N_{\mathrm{ns}}(A)
&=\frac{2\Delta}{\pi A}
\left[
2A\int_0^{\theta_0}\sin^2\theta\dd\theta
+\int_{\theta_0}^{\pi-\theta_0}\sin\theta\dd\theta
\right]\nonumber\\
&=\frac{2\Delta}{\pi A}
\left[
A\theta_0-\frac{A}{2}\sin(2\theta_0)
+2\cos\theta_0
\right]\nonumber\\
&=\frac{2\Delta}{\pi A}
\left[A\theta_0+\cos\theta_0\right].
\label{eq:df_sat_pasos}
\end{align}
En el último paso se usó
\[
\frac{A}{2}\sin(2\theta_0)
=A\sin\theta_0\cos\theta_0
=\cos\theta_0.
\]
Finalmente,
\[
\cos\theta_0
=\sqrt{1-\frac1{A^2}}
=\frac{\sqrt{A^2-1}}{A},
\]
y \cref{eq:df_sat_pasos} se convierte en
\[
N_{\mathrm{ns}}(A)
=\Delta\frac{2}{\pi}
\left[
\arcsin\left(\frac1A\right)
+\frac{\sqrt{A^2-1}}{A^2}
\right],
\]
que coincide con \cref{eq:df_saturacion}.

\fi
\ifnum\HAFOderivationpart=7\relax
\subsection{Derivación completa de la saturación sesgada}
\label{app:df_saturacion_sesgada}

Considérese
\[
u(\theta)=c+A\cos\theta,\qquad A>0.
\]
Para escribir una fórmula que incluya los casos parcialmente y completamente
saturados, se define
\[
[v]_{-1}^{1}=\max(-1,\min(1,v))
\]
y los ángulos
\begin{equation}
\theta_+
=\arccos\left[\frac{1-c}{A}\right]_{-1}^{1},
\qquad
\theta_-
=\arccos\left[\frac{-1-c}{A}\right]_{-1}^{1}.
\label{eq:angulos_saturacion_sesgada}
\end{equation}
Como \((1-c)/A>(-1-c)/A\) y \(\arccos\) es decreciente,
\(0\leq\theta_+\leq\theta_-\leq\pi\). En el intervalo
\([0,\pi]\),
\[
\operatorname{sat}(u(\theta))=
\begin{cases}
1,&0\leq\theta\leq\theta_+,\\
c+A\cos\theta,&\theta_+\leq\theta\leq\theta_-,\\
-1,&\theta_-\leq\theta\leq\pi.
\end{cases}
\]
La simetría respecto de \(2\pi\) permite integrar solamente en
\([0,\pi]\). Para \(\psi=\Delta\operatorname{sat}\),
\begin{align}
\Psi_0(c,A)
=\frac{\Delta}{\pi}\Big[
&\theta_+
+c(\theta_--\theta_+)
+A(\sin\theta_--\sin\theta_+)\nonumber\\
&-(\pi-\theta_-)
\Big].
\label{eq:psi0_saturacion_cerrada}
\end{align}

Para la componente fundamental se usan
\[
\int\cos\theta\dd\theta=\sin\theta,\qquad
\int\cos^2\theta\dd\theta
=\frac{\theta}{2}+\frac{\sin2\theta}{4}.
\]
La sustitución por intervalos da
\begin{align}
\Psi_1(c,A)
=\frac{2\Delta}{\pi}\Bigg[
&\sin\theta_++\sin\theta_-
+c(\sin\theta_--\sin\theta_+)\nonumber\\
&+\frac{A}{2}(\theta_--\theta_+)
+\frac{A}{4}
\bigl(\sin2\theta_--\sin2\theta_+\bigr)
\Bigg].
\label{eq:psi1_saturacion_cerrada}
\end{align}
Las operaciones de recorte en
\cref{eq:angulos_saturacion_sesgada} hacen que
\cref{eq:psi0_saturacion_cerrada,eq:psi1_saturacion_cerrada} sigan siendo
válidas cuando alguna de las tres regiones tiene longitud cero.

La ecuación de media puede reducirse a una igualdad escalar. De
\cref{eq:transferencia_cramer},
\begin{align}
W_q(0)
&=-\alpha
\frac{\beta+\gamma}
{\alpha(1+\ell)(\beta+\gamma)-\alpha\gamma}\nonumber\\
&=-\frac{\beta+\gamma}
{(1+\ell)(\beta+\gamma)-\gamma}.
\label{eq:ganancia_dc}
\end{align}
Por tanto,
\begin{equation}
c=W_q(0)\Psi_0(c,A),
\label{eq:cierre_dc_escalar}
\end{equation}
y el cierre armónico es
\[
1-W_q(\ii\omega)\frac{\Psi_1(c,A)}{A}=0.
\]
Estas tres ecuaciones reales ---una para la media, y las partes real e
imaginaria del cierre--- determinan \((c,A,\omega)\).

\fi
\ifnum\HAFOderivationpart=8\relax
\subsection{Derivación completa para la arctangente}
\label{app:df_arctan}

Sea
\[
\psi_{\mathrm a}(\sigma)=a_2\arctan(\rho\sigma),
\qquad
\xi=\rho A.
\]
Se define
\[
I(\xi)=
\int_0^\pi\arctan(\xi\sin\theta)\sin\theta\dd\theta.
\]
Como \(I(0)=0\), puede calcularse \(I\) integrando su derivada:
\begin{align}
I'(\xi)
&=\int_0^\pi
\frac{\sin^2\theta}{1+\xi^2\sin^2\theta}\dd\theta\nonumber\\
&=\frac1{\xi^2}
\left[
\pi-\int_0^\pi
\frac{\dd\theta}{1+\xi^2\sin^2\theta}
\right]\nonumber\\
&=\frac{\pi}{\xi^2}
\left(1-\frac1{\sqrt{1+\xi^2}}\right).
\label{eq:derivada_integral_arctan}
\end{align}
En la última igualdad se usó la integral estándar
\[
\int_0^\pi\frac{\dd\theta}{1+\xi^2\sin^2\theta}
=\frac{\pi}{\sqrt{1+\xi^2}}.
\]
Una primitiva que se anula en \(\xi=0\) es
\begin{equation}
I(\xi)
=\pi\frac{\sqrt{1+\xi^2}-1}{\xi}.
\label{eq:integral_arctan}
\end{equation}
Al sustituir \(\xi=\rho A\) en
\cref{eq:df_centrada},
\begin{align}
N_{\mathrm a}(A)
&=\frac{2a_2}{\pi A}I(\rho A)\nonumber\\
&=\frac{2a_2}{\rho A^2}
\left(\sqrt{1+\rho^2A^2}-1\right)\nonumber\\
&=\frac{2a_2\rho}
{1+\sqrt{1+\rho^2A^2}}.
\label{eq:df_arctan_derivada}
\end{align}

Si el cierre exige \(N_{\mathrm a}(A)=k\), la amplitud puede despejarse.
De la forma racionalizada,
\[
1+\sqrt{1+\rho^2A^2}=\frac{2a_2\rho}{k}.
\]
Al aislar la raíz, elevar al cuadrado y simplificar,
\begin{equation}
A^2
=\frac{4a_2(a_2\rho-k)}{k^2\rho}.
\label{eq:amplitud_arctan}
\end{equation}
La solución se acepta solamente si el lado derecho es positivo y la igualdad
previa a elevar al cuadrado tiene lado derecho no menor que \(2\). De manera
equivalente, para \(A>0\), \(k\) tiene el mismo signo que \(a_2\) y se
encuentra estrictamente entre \(0\) y \(a_2\rho\), con el orden de los
extremos ajustado al signo de \(a_2\).

\fi
\ifnum\HAFOderivationpart=9\relax
\subsection{Derivación completa de la identidad modal y su normalización}
\label{app:identidad_modal}

Sea \(P_0=P+bkr^{\T}\). Entonces
\[
\zeta I-P_0=(\zeta I-P)-bkr^{\T}.
\]
Si \(\zeta I-P\) es invertible, el lema del determinante matricial
\(\det(M+uv^{\T})=\det(M)(1+v^{\T}M^{-1}u)\), con
\[
M=\zeta I-P,\qquad u=-bk,\qquad v=r,
\]
produce
\begin{align}
\det(\zeta I-P_0)
&=\det(\zeta I-P)
\left[1-kr^{\T}(\zeta I-P)^{-1}b\right]\nonumber\\
&=\det(\zeta I-P)\bigl[1-kG(\zeta)\bigr],
\label{eq:lema_determinante_aplicado}
\end{align}
donde
\[
G(\zeta)=r^{\T}(\zeta I-P)^{-1}b,
\qquad W_q(s)=G(s^q).
\]
Por continuidad, la identidad polinómica se extiende también a los valores en
los que la resolvente de \(P\) es singular. Si
\[
1-kG(\zeta_0)=0,
\]
entonces \(\det(\zeta_0I-P_0)=0\), de modo que existe
\(v\neq0\) con \(P_0v=\zeta_0v\). Siempre que \(r^{\T}v\neq0\), se reemplaza
\(v\) por \(v/(r^{\T}v)\) y se obtiene \(r^{\T}v=1\). Ésta es la
normalización usada en \cref{eq:modo_normalizado,eq:semilla_modal}.

En la rama sesgada,
\[
X_1=\bigl[(\ii\omega_0)^qI-P\bigr]^{-1}b\Psi_1.
\]
Al premultiplicar por \(r^{\T}\),
\[
r^{\T}X_1
=W_q(\ii\omega_0)\Psi_1
=W_q(\ii\omega_0)N_1A.
\]
El cierre \(1-W_qN_1=0\) prueba algebraicamente que
\(r^{\T}X_1=A\).

\fi
\ifnum\HAFOderivationpart=10\relax
\subsection{Derivación completa de la homotopía afín}
\label{app:homotopia_afin}

La aproximación sesgada separa
\[
\psi(\sigma)
=\Psi_0+N_1(\sigma-c)
+\bigl[\psi(\sigma)-\Psi_0-N_1(\sigma-c)\bigr].
\]
Se definen
\begin{equation}
P_{\mathrm a}=P+N_1br^{\T},\qquad
d_{\mathrm a}=b(\Psi_0-N_1c).
\label{eq:parametros_homotopia_afin}
\end{equation}
La familia de campos es
\begin{equation}
\begin{aligned}
F_\varepsilon(X)
&=P_{\mathrm a}X+d_{\mathrm a}\\
&\quad+\varepsilon b
\left[
\psi(r^{\T}X)-\Psi_0-N_1(r^{\T}X-c)
\right],
\qquad 0\leq\varepsilon\leq1.
\end{aligned}
\label{eq:homotopia_afin}
\end{equation}

En \(\varepsilon=0\), la media \(\bar X\) de
\cref{eq:cierre_dc_sesgado} es un equilibrio del auxiliar:
\begin{align*}
F_0(\bar X)
&=P\bar X+bN_1c+b\Psi_0-bN_1c\\
&=P\bar X+b\Psi_0=0.
\end{align*}
Para una perturbación \(\xi=X-\bar X\),
\[
F_0(\bar X+\xi)=P_{\mathrm a}\xi,
\]
de modo que el auxiliar conserva tanto la media como el modo armónico.

En \(\varepsilon=1\), al agrupar términos,
\begin{align*}
F_1(X)
&=PX+bN_1r^{\T}X+b\Psi_0-bN_1c\\
&\quad+b\psi(r^{\T}X)-b\Psi_0-bN_1r^{\T}X+bN_1c\\
&=PX+b\psi(r^{\T}X).
\end{align*}
Así, el extremo final de \cref{eq:homotopia_afin} es exactamente el sistema
físico de \cref{eq:lure_comun}. Una homotopía centrada de la forma
\(P_0X+\varepsilon b[\psi(\sigma)-k\sigma]\) no conserva, en general, la
media \(c\neq0\); por eso no se utiliza para la rama sesgada.
\fi

\endgroup
La comprobación fija también el signo: el vector \(b\) contiene \(-\alpha\), mientras que
\(\psi\) conserva el signo de sus parámetros. Esta convención determina el
signo empleado en el balance armónico.

En \cref{eq:lure_comun}, \(q\) designa únicamente el orden fraccionario. El
vector de entrada se denota por \(b\); así, ambos símbolos designan objetos
matemáticos distintos. Para \(q=1\), el miembro izquierdo se reemplaza por
\(\dot X\) y se recupera la representación de Lur'e de orden entero utilizada
por Leonov y Kuznetsov
\cite{LeonovKuznetsov2013,KuznetsovEtAl2017}.

\subsection{Transferencia obtenida de la transformada de Laplace}

Sea
\[
U(s)=\mathcal L_u\{\psi(\sigma(t_0+u))\}(s),\qquad
\widehat X(s)=\mathcal L_u\{X(t_0+u)\}(s).
\]
Al aplicar \cref{eq:laplace_caputo} a \cref{eq:lure_comun} resulta
\begin{equation}
(s^qI-P)\widehat X(s)
=s^{q-1}X(t_0)+b\,U(s).
\label{eq:lure_laplace}
\end{equation}
Por tanto, la salida \(\Sigma(s)=r^{\T}\widehat X(s)\) es
\begin{equation}
\begin{aligned}
\Sigma(s)
&=r^{\T}(s^qI-P)^{-1}s^{q-1}X(t_0)
  +W_q(s)U(s),\\
W_q(s)&=r^{\T}(s^qI-P)^{-1}b.
\end{aligned}
\label{eq:transferencia_fraccionaria}
\end{equation}
El primer término de \cref{eq:transferencia_fraccionaria} es la respuesta
debida a la condición inicial. La transferencia \(W_q\) es, por definición, la
relación entre la entrada y la salida cuando se separa ese término. Se obtiene
directamente al aplicar la transformada de Laplace a la derivada de Caputo y
resolver el bloque lineal.

El numerador y el denominador se derivan mediante menores, expansión del
determinante y regla de Cramer.
\begingroup
\def\HAFOderivationpart{2}
% Este archivo es una biblioteca de derivaciones. Cada bloque se inserta en el
% lugar del cuerpo donde se usa, mediante \HAFOderivationpart. No se imprime
% como apéndice independiente.
\ifnum\HAFOderivationpart=1\relax
\subsection{Derivación completa de la forma de Lur'e}
\label{app:verificacion_lure}

Sea \(h(x)=\ell x+\psi(x)\). La primera ecuación de los dos modelos es
\[
{}^{\mathrm C}D_{t_0}^{q}x
=\alpha(y-x-h(x)).
\]
Al distribuir \(\alpha\),
\begin{align}
\alpha(y-x-h(x))
&=\alpha y-\alpha x-\alpha\ell x-\alpha\psi(x)\nonumber\\
&=-\alpha(1+\ell)x+\alpha y-\alpha\psi(x).
\label{eq:expansion_primera_fila}
\end{align}
Por otra parte,
\[
PX=
\begin{pmatrix}
-\alpha(1+\ell)x+\alpha y\\
x-y+z\\
-\beta y-\gamma z
\end{pmatrix},
\qquad
b\psi(r^{\T}X)=
\begin{pmatrix}
-\alpha\psi(x)\\0\\0
\end{pmatrix}.
\]
La suma reproduce \cref{eq:expansion_primera_fila} y las otras dos
ecuaciones. Para la saturación se toma
\(\ell=m_1\) y
\(\psi=(m_0-m_1)\operatorname{sat}\); para la arctangente,
\(\ell=a_1\) y \(\psi=a_2\arctan(\rho\,\cdot)\).

\fi
\ifnum\HAFOderivationpart=2\relax
\subsection{Derivación completa del determinante y la transferencia}
\label{app:transferencia_explicita}

Con \(\lambda=s^q\),
\begin{equation}
\lambda I-P=
\begin{pmatrix}
\lambda+\alpha(1+\ell)&-\alpha&0\\
-1&\lambda+1&-1\\
0&\beta&\lambda+\gamma
\end{pmatrix}.
\label{eq:matriz_resolvente}
\end{equation}
El menor de la posición \((1,1)\) es
\begin{align}
Q(\lambda)
&=
\det
\begin{pmatrix}
\lambda+1&-1\\
\beta&\lambda+\gamma
\end{pmatrix}\nonumber\\
&=(\lambda+1)(\lambda+\gamma)+\beta.
\label{eq:menor_Q}
\end{align}
Al desarrollar \(\det(\lambda I-P)\) por la primera fila,
\begin{align}
\Delta_\ell(\lambda)
&=\bigl[\lambda+\alpha(1+\ell)\bigr]Q(\lambda)
-(-\alpha)
\det
\begin{pmatrix}
-1&-1\\
0&\lambda+\gamma
\end{pmatrix}\nonumber\\
&=\bigl[\lambda+\alpha(1+\ell)\bigr]Q(\lambda)
-\alpha(\lambda+\gamma).
\label{eq:determinante_transferencia}
\end{align}

Para calcular la primera componente de
\((\lambda I-P)^{-1}b\), la regla de Cramer sustituye la primera columna de
\cref{eq:matriz_resolvente} por
\(b=(-\alpha,0,0)^{\T}\). El determinante del numerador es
\[
\det
\begin{pmatrix}
-\alpha&-\alpha&0\\
0&\lambda+1&-1\\
0&\beta&\lambda+\gamma
\end{pmatrix}
=-\alpha Q(\lambda).
\]
Como \(r^{\T}\) selecciona esa primera componente,
\begin{equation}
W_q(s)
=r^{\T}(s^qI-P)^{-1}b
=-\alpha\frac{Q(s^q)}{\Delta_\ell(s^q)}.
\label{eq:transferencia_cramer}
\end{equation}

La identidad
\[
P-s^qI=-(s^qI-P)
\]
implica
\[
r^{\T}(P-s^qI)^{-1}b=-W_q(s).
\]
Por ello, cambiar el orden de los términos dentro de la resolvente cambia
simultáneamente el signo de la transferencia y el signo del cierre. No cambia
la solución física si ambas modificaciones se realizan de manera coherente.

\fi
\ifnum\HAFOderivationpart=3\relax
\subsection{Derivación completa de los equilibrios}
\label{app:equilibrios}

En un equilibrio \(e=(x_e,y_e,z_e)^{\T}\), las dos últimas ecuaciones de
ambos sistemas dan
\[
0=x_e-y_e+z_e,\qquad
0=-\beta y_e-\gamma z_e.
\]
De la primera,
\[
z_e=y_e-x_e.
\]
Al sustituirla en la segunda,
\begin{align*}
0
&=-\beta y_e-\gamma(y_e-x_e)\\
&=-(\beta+\gamma)y_e+\gamma x_e.
\end{align*}
Si \(\beta+\gamma\neq0\),
\begin{equation}
y_e=\frac{\gamma}{\beta+\gamma}x_e,\qquad
z_e=-\frac{\beta}{\beta+\gamma}x_e.
\label{eq:equilibrios_yz}
\end{equation}
La primera ecuación de equilibrio exige
\[
0=y_e-x_e-h(x_e).
\]
Mediante \cref{eq:equilibrios_yz},
\begin{equation}
h(x_e)+\frac{\beta}{\beta+\gamma}x_e=0.
\label{eq:equilibrio_escalar_comun}
\end{equation}
Esta única ecuación escalar determina todos los equilibrios; las otras dos
coordenadas se recuperan después con \cref{eq:equilibrios_yz}.

\subsubsection{Equilibrios del Chua no suave}

En la región \(|x_e|<1\), \(h(x_e)=m_0x_e\). Entonces
\begin{equation}
\left(m_0+\frac{\beta}{\beta+\gamma}\right)x_e=0.
\label{eq:equilibrio_interior_ns}
\end{equation}
Cuando el coeficiente no es cero, la única raíz interior es \(x_e=0\).

Para \(x_e>1\),
\[
h(x_e)=m_1x_e+(m_0-m_1),
\]
y \cref{eq:equilibrio_escalar_comun} queda
\[
\left(m_1+\frac{\beta}{\beta+\gamma}\right)x_e+(m_0-m_1)=0.
\]
Por tanto,
\begin{equation}
x_\star=
\frac{m_1-m_0}
{m_1+\dfrac{\beta}{\beta+\gamma}}.
\label{eq:raiz_exterior_ns}
\end{equation}
Si \(x_\star>1\), existe el equilibrio positivo. Como la no linealidad es
impar, existe también la raíz \(-x_\star<-1\). Los tres equilibrios son
\begin{equation}
E_0=(0,0,0),\qquad
E_\pm=
\left(
\pm x_\star,\,
\pm\frac{\gamma}{\beta+\gamma}x_\star,\,
\mp\frac{\beta}{\beta+\gamma}x_\star
\right).
\label{eq:equilibrios_ns_simbolicos}
\end{equation}
La condición regional \(x_\star>1\) debe comprobarse; no basta con resolver la
ecuación lineal exterior.

\subsubsection{Equilibrios del Chua con arctangente}

Al usar
\(h(x)=a_1x+a_2\arctan(\rho x)\) en
\cref{eq:equilibrio_escalar_comun}, se obtiene
\begin{equation}
\left(a_1+\frac{\beta}{\beta+\gamma}\right)x_e
+a_2\arctan(\rho x_e)=0.
\label{eq:equilibrio_escalar_arctan}
\end{equation}
La expresión es impar, de modo que \(x_e=0\) siempre es raíz y toda raíz
positiva no nula tiene una compañera negativa. Las raíces no nulas se
determinan numéricamente y se sustituyen en \cref{eq:equilibrios_yz}. Esta
ecuación se evalúa con \(\rho=1\) para el caso bibliográfico y con el valor de
\(\rho\) de \cref{eq:parametros_c590} para el caso \texttt{c590}; no deben
intercambiarse las dos parametrizaciones.

\fi
\ifnum\HAFOderivationpart=4\relax
\subsection{Derivación completa de los jacobianos regionales}
\label{app:jacobianos}

Para la forma de Lur'e,
\[
F(X)=PX+b\psi(r^{\T}X).
\]
La regla de la cadena produce
\begin{equation}
J_e=P+b\,\psi'(r^{\T}e)r^{\T}.
\label{eq:jacobiano_lure}
\end{equation}
Como \(r^{\T}e=x_e\), solamente cambia la entrada \((1,1)\).

En el sistema no suave,
\[
\psi'_{\mathrm{ns}}(x)=
\begin{cases}
m_0-m_1,&|x|<1,\\
0,&|x|>1.
\end{cases}
\]
Por ello,
\begin{equation}
J_{\mathrm{in}}=
\begin{pmatrix}
-\alpha(1+m_0)&\alpha&0\\
1&-1&1\\
0&-\beta&-\gamma
\end{pmatrix},
\qquad
J_{\mathrm{out}}=
\begin{pmatrix}
-\alpha(1+m_1)&\alpha&0\\
1&-1&1\\
0&-\beta&-\gamma
\end{pmatrix}.
\label{eq:jacobianos_ns}
\end{equation}
En \(x=\pm1\) la derivada clásica no existe. Los equilibrios usados en este
reporte no se encuentran sobre esas superficies, por lo que se emplea el
jacobiano de la región que contiene a cada uno.

Para la arctangente,
\[
\psi'_{\mathrm a}(x)
=\frac{a_2\rho}{1+\rho^2x^2},
\]
y
\begin{equation}
J_e=
\begin{pmatrix}
\displaystyle
-\alpha\left(1+a_1+
\frac{a_2\rho}{1+\rho^2x_e^2}\right)
&\alpha&0\\[3mm]
1&-1&1\\
0&-\beta&-\gamma
\end{pmatrix}.
\label{eq:jacobiano_arctan}
\end{equation}
Los valores propios de
\cref{eq:jacobianos_ns,eq:jacobiano_arctan} se clasifican con
\cref{eq:matignon}; la estabilidad local no sustituye el sondeo de cuenca.

\fi
\ifnum\HAFOderivationpart=5\relax
\subsection{Derivación completa del polinomio característico}
\label{app:polinomio_caracteristico}

Sea \(g=h'(x_e)\) la pendiente total de la no linealidad en un equilibrio.
El jacobiano de cualquiera de las dos familias puede escribirse como
\begin{equation}
J(g)=
\begin{pmatrix}
-\alpha(1+g)&\alpha&0\\
1&-1&1\\
0&-\beta&-\gamma
\end{pmatrix}.
\label{eq:jacobiano_pendiente_total}
\end{equation}
Por tanto,
\[
\lambda I-J(g)=
\begin{pmatrix}
\lambda+\alpha(1+g)&-\alpha&0\\
-1&\lambda+1&-1\\
0&\beta&\lambda+\gamma
\end{pmatrix}.
\]
El desarrollo por la primera fila, igual al de
\cref{eq:determinante_transferencia}, produce
\begin{equation}
\chi_g(\lambda)
=\bigl[\lambda+\alpha(1+g)\bigr]
\bigl[(\lambda+1)(\lambda+\gamma)+\beta\bigr]
-\alpha(\lambda+\gamma).
\label{eq:polinomio_caracteristico_agrupado}
\end{equation}
Como
\[
(\lambda+1)(\lambda+\gamma)+\beta
=\lambda^2+(1+\gamma)\lambda+(\beta+\gamma),
\]
la multiplicación y agrupación por potencias de \(\lambda\) da
\begin{equation}
\begin{aligned}
\chi_g(\lambda)
={}&\lambda^3+
\bigl[1+\gamma+\alpha(1+g)\bigr]\lambda^2\\
&+\bigl[
\beta+\gamma+\alpha(1+g)(1+\gamma)-\alpha
\bigr]\lambda\\
&+\alpha\bigl[(1+g)(\beta+\gamma)-\gamma\bigr].
\end{aligned}
\label{eq:polinomio_caracteristico_expandido}
\end{equation}
En el sistema no suave se sustituye \(g=m_0\) en \(E_0\) y \(g=m_1\)
en \(E_\pm\). En el sistema con arctangente,
\[
g=a_1+\frac{a_2\rho}{1+\rho^2x_e^2}.
\]
Las raíces de \cref{eq:polinomio_caracteristico_expandido} son los valores
propios utilizados en \cref{eq:matignon}; así, la clasificación local se
obtiene a partir del mismo campo que define la integración.

\fi
\ifnum\HAFOderivationpart=6\relax
\subsection{Derivación completa para la saturación centrada}
\label{app:df_saturacion_centrada}

Sea
\[
\psi_{\mathrm{ns}}(\sigma)=\Delta\operatorname{sat}(\sigma),
\qquad \Delta=m_0-m_1.
\]
Si \(0<A\leq1\), no hay saturación y
\[
\psi_{\mathrm{ns}}(A\sin\theta)=\Delta A\sin\theta.
\]
Como \(\int_0^\pi\sin^2\theta\dd\theta=\pi/2\),
\[
N_{\mathrm{ns}}(A)
=\frac{2}{\pi A}\Delta A\frac{\pi}{2}
=\Delta.
\]

Para \(A>1\), defínase
\[
\theta_0=\arcsin\left(\frac1A\right),
\qquad 0<\theta_0<\frac{\pi}{2}.
\]
En \(0\leq\theta\leq\theta_0\) y
\(\pi-\theta_0\leq\theta\leq\pi\), la entrada permanece en la parte lineal;
en el intervalo central la saturación vale \(1\). Entonces
\begin{align}
N_{\mathrm{ns}}(A)
&=\frac{2\Delta}{\pi A}
\left[
2A\int_0^{\theta_0}\sin^2\theta\dd\theta
+\int_{\theta_0}^{\pi-\theta_0}\sin\theta\dd\theta
\right]\nonumber\\
&=\frac{2\Delta}{\pi A}
\left[
A\theta_0-\frac{A}{2}\sin(2\theta_0)
+2\cos\theta_0
\right]\nonumber\\
&=\frac{2\Delta}{\pi A}
\left[A\theta_0+\cos\theta_0\right].
\label{eq:df_sat_pasos}
\end{align}
En el último paso se usó
\[
\frac{A}{2}\sin(2\theta_0)
=A\sin\theta_0\cos\theta_0
=\cos\theta_0.
\]
Finalmente,
\[
\cos\theta_0
=\sqrt{1-\frac1{A^2}}
=\frac{\sqrt{A^2-1}}{A},
\]
y \cref{eq:df_sat_pasos} se convierte en
\[
N_{\mathrm{ns}}(A)
=\Delta\frac{2}{\pi}
\left[
\arcsin\left(\frac1A\right)
+\frac{\sqrt{A^2-1}}{A^2}
\right],
\]
que coincide con \cref{eq:df_saturacion}.

\fi
\ifnum\HAFOderivationpart=7\relax
\subsection{Derivación completa de la saturación sesgada}
\label{app:df_saturacion_sesgada}

Considérese
\[
u(\theta)=c+A\cos\theta,\qquad A>0.
\]
Para escribir una fórmula que incluya los casos parcialmente y completamente
saturados, se define
\[
[v]_{-1}^{1}=\max(-1,\min(1,v))
\]
y los ángulos
\begin{equation}
\theta_+
=\arccos\left[\frac{1-c}{A}\right]_{-1}^{1},
\qquad
\theta_-
=\arccos\left[\frac{-1-c}{A}\right]_{-1}^{1}.
\label{eq:angulos_saturacion_sesgada}
\end{equation}
Como \((1-c)/A>(-1-c)/A\) y \(\arccos\) es decreciente,
\(0\leq\theta_+\leq\theta_-\leq\pi\). En el intervalo
\([0,\pi]\),
\[
\operatorname{sat}(u(\theta))=
\begin{cases}
1,&0\leq\theta\leq\theta_+,\\
c+A\cos\theta,&\theta_+\leq\theta\leq\theta_-,\\
-1,&\theta_-\leq\theta\leq\pi.
\end{cases}
\]
La simetría respecto de \(2\pi\) permite integrar solamente en
\([0,\pi]\). Para \(\psi=\Delta\operatorname{sat}\),
\begin{align}
\Psi_0(c,A)
=\frac{\Delta}{\pi}\Big[
&\theta_+
+c(\theta_--\theta_+)
+A(\sin\theta_--\sin\theta_+)\nonumber\\
&-(\pi-\theta_-)
\Big].
\label{eq:psi0_saturacion_cerrada}
\end{align}

Para la componente fundamental se usan
\[
\int\cos\theta\dd\theta=\sin\theta,\qquad
\int\cos^2\theta\dd\theta
=\frac{\theta}{2}+\frac{\sin2\theta}{4}.
\]
La sustitución por intervalos da
\begin{align}
\Psi_1(c,A)
=\frac{2\Delta}{\pi}\Bigg[
&\sin\theta_++\sin\theta_-
+c(\sin\theta_--\sin\theta_+)\nonumber\\
&+\frac{A}{2}(\theta_--\theta_+)
+\frac{A}{4}
\bigl(\sin2\theta_--\sin2\theta_+\bigr)
\Bigg].
\label{eq:psi1_saturacion_cerrada}
\end{align}
Las operaciones de recorte en
\cref{eq:angulos_saturacion_sesgada} hacen que
\cref{eq:psi0_saturacion_cerrada,eq:psi1_saturacion_cerrada} sigan siendo
válidas cuando alguna de las tres regiones tiene longitud cero.

La ecuación de media puede reducirse a una igualdad escalar. De
\cref{eq:transferencia_cramer},
\begin{align}
W_q(0)
&=-\alpha
\frac{\beta+\gamma}
{\alpha(1+\ell)(\beta+\gamma)-\alpha\gamma}\nonumber\\
&=-\frac{\beta+\gamma}
{(1+\ell)(\beta+\gamma)-\gamma}.
\label{eq:ganancia_dc}
\end{align}
Por tanto,
\begin{equation}
c=W_q(0)\Psi_0(c,A),
\label{eq:cierre_dc_escalar}
\end{equation}
y el cierre armónico es
\[
1-W_q(\ii\omega)\frac{\Psi_1(c,A)}{A}=0.
\]
Estas tres ecuaciones reales ---una para la media, y las partes real e
imaginaria del cierre--- determinan \((c,A,\omega)\).

\fi
\ifnum\HAFOderivationpart=8\relax
\subsection{Derivación completa para la arctangente}
\label{app:df_arctan}

Sea
\[
\psi_{\mathrm a}(\sigma)=a_2\arctan(\rho\sigma),
\qquad
\xi=\rho A.
\]
Se define
\[
I(\xi)=
\int_0^\pi\arctan(\xi\sin\theta)\sin\theta\dd\theta.
\]
Como \(I(0)=0\), puede calcularse \(I\) integrando su derivada:
\begin{align}
I'(\xi)
&=\int_0^\pi
\frac{\sin^2\theta}{1+\xi^2\sin^2\theta}\dd\theta\nonumber\\
&=\frac1{\xi^2}
\left[
\pi-\int_0^\pi
\frac{\dd\theta}{1+\xi^2\sin^2\theta}
\right]\nonumber\\
&=\frac{\pi}{\xi^2}
\left(1-\frac1{\sqrt{1+\xi^2}}\right).
\label{eq:derivada_integral_arctan}
\end{align}
En la última igualdad se usó la integral estándar
\[
\int_0^\pi\frac{\dd\theta}{1+\xi^2\sin^2\theta}
=\frac{\pi}{\sqrt{1+\xi^2}}.
\]
Una primitiva que se anula en \(\xi=0\) es
\begin{equation}
I(\xi)
=\pi\frac{\sqrt{1+\xi^2}-1}{\xi}.
\label{eq:integral_arctan}
\end{equation}
Al sustituir \(\xi=\rho A\) en
\cref{eq:df_centrada},
\begin{align}
N_{\mathrm a}(A)
&=\frac{2a_2}{\pi A}I(\rho A)\nonumber\\
&=\frac{2a_2}{\rho A^2}
\left(\sqrt{1+\rho^2A^2}-1\right)\nonumber\\
&=\frac{2a_2\rho}
{1+\sqrt{1+\rho^2A^2}}.
\label{eq:df_arctan_derivada}
\end{align}

Si el cierre exige \(N_{\mathrm a}(A)=k\), la amplitud puede despejarse.
De la forma racionalizada,
\[
1+\sqrt{1+\rho^2A^2}=\frac{2a_2\rho}{k}.
\]
Al aislar la raíz, elevar al cuadrado y simplificar,
\begin{equation}
A^2
=\frac{4a_2(a_2\rho-k)}{k^2\rho}.
\label{eq:amplitud_arctan}
\end{equation}
La solución se acepta solamente si el lado derecho es positivo y la igualdad
previa a elevar al cuadrado tiene lado derecho no menor que \(2\). De manera
equivalente, para \(A>0\), \(k\) tiene el mismo signo que \(a_2\) y se
encuentra estrictamente entre \(0\) y \(a_2\rho\), con el orden de los
extremos ajustado al signo de \(a_2\).

\fi
\ifnum\HAFOderivationpart=9\relax
\subsection{Derivación completa de la identidad modal y su normalización}
\label{app:identidad_modal}

Sea \(P_0=P+bkr^{\T}\). Entonces
\[
\zeta I-P_0=(\zeta I-P)-bkr^{\T}.
\]
Si \(\zeta I-P\) es invertible, el lema del determinante matricial
\(\det(M+uv^{\T})=\det(M)(1+v^{\T}M^{-1}u)\), con
\[
M=\zeta I-P,\qquad u=-bk,\qquad v=r,
\]
produce
\begin{align}
\det(\zeta I-P_0)
&=\det(\zeta I-P)
\left[1-kr^{\T}(\zeta I-P)^{-1}b\right]\nonumber\\
&=\det(\zeta I-P)\bigl[1-kG(\zeta)\bigr],
\label{eq:lema_determinante_aplicado}
\end{align}
donde
\[
G(\zeta)=r^{\T}(\zeta I-P)^{-1}b,
\qquad W_q(s)=G(s^q).
\]
Por continuidad, la identidad polinómica se extiende también a los valores en
los que la resolvente de \(P\) es singular. Si
\[
1-kG(\zeta_0)=0,
\]
entonces \(\det(\zeta_0I-P_0)=0\), de modo que existe
\(v\neq0\) con \(P_0v=\zeta_0v\). Siempre que \(r^{\T}v\neq0\), se reemplaza
\(v\) por \(v/(r^{\T}v)\) y se obtiene \(r^{\T}v=1\). Ésta es la
normalización usada en \cref{eq:modo_normalizado,eq:semilla_modal}.

En la rama sesgada,
\[
X_1=\bigl[(\ii\omega_0)^qI-P\bigr]^{-1}b\Psi_1.
\]
Al premultiplicar por \(r^{\T}\),
\[
r^{\T}X_1
=W_q(\ii\omega_0)\Psi_1
=W_q(\ii\omega_0)N_1A.
\]
El cierre \(1-W_qN_1=0\) prueba algebraicamente que
\(r^{\T}X_1=A\).

\fi
\ifnum\HAFOderivationpart=10\relax
\subsection{Derivación completa de la homotopía afín}
\label{app:homotopia_afin}

La aproximación sesgada separa
\[
\psi(\sigma)
=\Psi_0+N_1(\sigma-c)
+\bigl[\psi(\sigma)-\Psi_0-N_1(\sigma-c)\bigr].
\]
Se definen
\begin{equation}
P_{\mathrm a}=P+N_1br^{\T},\qquad
d_{\mathrm a}=b(\Psi_0-N_1c).
\label{eq:parametros_homotopia_afin}
\end{equation}
La familia de campos es
\begin{equation}
\begin{aligned}
F_\varepsilon(X)
&=P_{\mathrm a}X+d_{\mathrm a}\\
&\quad+\varepsilon b
\left[
\psi(r^{\T}X)-\Psi_0-N_1(r^{\T}X-c)
\right],
\qquad 0\leq\varepsilon\leq1.
\end{aligned}
\label{eq:homotopia_afin}
\end{equation}

En \(\varepsilon=0\), la media \(\bar X\) de
\cref{eq:cierre_dc_sesgado} es un equilibrio del auxiliar:
\begin{align*}
F_0(\bar X)
&=P\bar X+bN_1c+b\Psi_0-bN_1c\\
&=P\bar X+b\Psi_0=0.
\end{align*}
Para una perturbación \(\xi=X-\bar X\),
\[
F_0(\bar X+\xi)=P_{\mathrm a}\xi,
\]
de modo que el auxiliar conserva tanto la media como el modo armónico.

En \(\varepsilon=1\), al agrupar términos,
\begin{align*}
F_1(X)
&=PX+bN_1r^{\T}X+b\Psi_0-bN_1c\\
&\quad+b\psi(r^{\T}X)-b\Psi_0-bN_1r^{\T}X+bN_1c\\
&=PX+b\psi(r^{\T}X).
\end{align*}
Así, el extremo final de \cref{eq:homotopia_afin} es exactamente el sistema
físico de \cref{eq:lure_comun}. Una homotopía centrada de la forma
\(P_0X+\varepsilon b[\psi(\sigma)-k\sigma]\) no conserva, en general, la
media \(c\neq0\); por eso no se utiliza para la rama sesgada.
\fi

\endgroup
Para reutilizar el resultado, se denota
\[
G(\lambda)=r^{\T}(\lambda I-P)^{-1}b.
\]
El cálculo anterior se resume como
\begin{equation}
\begin{aligned}
\Delta_\ell(\lambda)
&=\det(\lambda I-P)\\
&=\bigl[\lambda+\alpha(1+\ell)\bigr]Q(\lambda)
  -\alpha(\lambda+\gamma),\\
G(\lambda)&=-\alpha\frac{Q(\lambda)}{\Delta_\ell(\lambda)},\\
W_q(s)&=G(s^q).
\end{aligned}
\label{eq:transferencia_explicita}
\end{equation}
Sobre el eje imaginario se usa
\(\lambda=(\ii\omega)^q\) con la rama principal de
\cref{eq:rama_principal}.

\subsection{Convención de signo y condición de cierre}

Supóngase que la componente fundamental de la no linealidad se aproxima por
\[
\psi(\sigma(t))\simeq N\,\sigma(t).
\]
En el predictor estacionario de Liouville--Weyl, sea \(S_1\) el fasor de
la salida \(\sigma\). La transferencia del bloque lineal da
\[
S_1=W_q(\ii\omega)N S_1.
\]
Esta igualdad relaciona fasores estacionarios; la respuesta al dato inicial
de \cref{eq:transferencia_fraccionaria} se trata por separado en el PVI
causal. Una respuesta armónica no nula del predictor requiere
\begin{equation}
R(\omega,N)
=1-W_q(\ii\omega)N=0.
\label{eq:cierre_armonico}
\end{equation}
Cuando \(N\) es real, \cref{eq:cierre_armonico} equivale a
\begin{equation}
\Im W_q(\ii\omega_0)=0,\qquad
N=\frac{1}{\Re W_q(\ii\omega_0)}.
\label{eq:condiciones_fase_ganancia}
\end{equation}
Se exige además que la transferencia sea finita, que su parte real no sea
cero y que la ganancia resultante pertenezca al rango de la función
descriptiva. La condición de fase por sí sola no garantiza una amplitud
admisible.

Algunas referencias y rutinas escriben la transferencia como
\[
W_q^{-}(s)=r^{\T}(P-s^qI)^{-1}b=-W_q(s).
\]
Con esa definición, el mismo cierre se escribe
\[
1+W_q^{-}(\ii\omega)N=0,\qquad
N=-\frac{1}{\Re W_q^{-}(\ii\omega)}.
\]
Las dos escrituras son algebraicamente equivalentes. Se adopta la convención de
\cref{eq:transferencia_fraccionaria,eq:cierre_armonico}, que fija de manera
unívoca el signo del cierre.

Una similitud real \(X=SY\), con \(S\) invertible, transforma los datos
en \(\widetilde P=S^{-1}PS\), \(\widetilde b=S^{-1}b\) y
\(\widetilde r^{\T}=r^{\T}S\). La identidad
\[
 (s^qI-\widetilde P)^{-1}=S^{-1}(s^qI-P)^{-1}S
\]
da \(\widetilde W_q=W_q\). El modo y la semilla se transforman con
\(S^{-1}\), sin cambiar el cierre. Si se usa además un desplazamiento
\(X=SY+d\), debe conservarse el término afín \(S^{-1}Pd\) y el
sesgo \(r^{\T}d\) dentro de la no linealidad. Cuando \(d\) es un
equilibrio, \(Pd+b\psi(r^{\T}d)=0\) permite escribir exactamente
la característica desplazada
\(\psi(r^{\T}SY+r^{\T}d)-\psi(r^{\T}d)\).

\subsection{Función descriptiva centrada}

La convención de primer armónico sigue el
tratamiento clásico de la función descriptiva y sus extensiones para
no linealidades estáticas\cite{GelbVanderVelde1968,Atherton1975,AthertonTanYerogluKavuranYuce2014}.
Para una entrada \(\sigma=A\sin\theta\), \(A>0\), y una no linealidad
estática e impar, la proyección sobre el primer armónico es
\begin{equation}
N_\psi(A)
=\frac{2}{\pi A}\int_0^\pi
\psi(A\sin\theta)\sin\theta\dd\theta.
\label{eq:df_centrada}
\end{equation}
La integral de \cref{eq:df_centrada} no contiene un operador temporal. Por
ello, \(N_\psi\) depende de la no linealidad y de la amplitud, mientras que el
orden \(q\) entra en la transferencia mediante
\((\ii\omega)^q\). El carácter fraccionario entra mediante el bloque dinámico;
la integral de la función descriptiva conserva su definición estática.
La imparidad anula la media; la simetría de
\(\psi(A\sin\theta)\) entre los semiperiodos y su proyección ortogonal
producen el factor \(2/(\pi A)\). No se aplica una derivada temporal a
esta integral, de modo que sustituir \(N_\psi\) por \(N_\psi^q\) sería
una modificación adicional del cierre, no una consecuencia de Caputo.

Para el Chua no suave, sea
\[
\Delta=m_0-m_1,\qquad
\psi_{\mathrm{ns}}(\sigma)=\Delta\operatorname{sat}(\sigma).
\]
La evaluación exacta de \cref{eq:df_centrada} se obtiene por integración por
tramos:
\begingroup
\def\HAFOderivationpart{6}
% Este archivo es una biblioteca de derivaciones. Cada bloque se inserta en el
% lugar del cuerpo donde se usa, mediante \HAFOderivationpart. No se imprime
% como apéndice independiente.
\ifnum\HAFOderivationpart=1\relax
\subsection{Derivación completa de la forma de Lur'e}
\label{app:verificacion_lure}

Sea \(h(x)=\ell x+\psi(x)\). La primera ecuación de los dos modelos es
\[
{}^{\mathrm C}D_{t_0}^{q}x
=\alpha(y-x-h(x)).
\]
Al distribuir \(\alpha\),
\begin{align}
\alpha(y-x-h(x))
&=\alpha y-\alpha x-\alpha\ell x-\alpha\psi(x)\nonumber\\
&=-\alpha(1+\ell)x+\alpha y-\alpha\psi(x).
\label{eq:expansion_primera_fila}
\end{align}
Por otra parte,
\[
PX=
\begin{pmatrix}
-\alpha(1+\ell)x+\alpha y\\
x-y+z\\
-\beta y-\gamma z
\end{pmatrix},
\qquad
b\psi(r^{\T}X)=
\begin{pmatrix}
-\alpha\psi(x)\\0\\0
\end{pmatrix}.
\]
La suma reproduce \cref{eq:expansion_primera_fila} y las otras dos
ecuaciones. Para la saturación se toma
\(\ell=m_1\) y
\(\psi=(m_0-m_1)\operatorname{sat}\); para la arctangente,
\(\ell=a_1\) y \(\psi=a_2\arctan(\rho\,\cdot)\).

\fi
\ifnum\HAFOderivationpart=2\relax
\subsection{Derivación completa del determinante y la transferencia}
\label{app:transferencia_explicita}

Con \(\lambda=s^q\),
\begin{equation}
\lambda I-P=
\begin{pmatrix}
\lambda+\alpha(1+\ell)&-\alpha&0\\
-1&\lambda+1&-1\\
0&\beta&\lambda+\gamma
\end{pmatrix}.
\label{eq:matriz_resolvente}
\end{equation}
El menor de la posición \((1,1)\) es
\begin{align}
Q(\lambda)
&=
\det
\begin{pmatrix}
\lambda+1&-1\\
\beta&\lambda+\gamma
\end{pmatrix}\nonumber\\
&=(\lambda+1)(\lambda+\gamma)+\beta.
\label{eq:menor_Q}
\end{align}
Al desarrollar \(\det(\lambda I-P)\) por la primera fila,
\begin{align}
\Delta_\ell(\lambda)
&=\bigl[\lambda+\alpha(1+\ell)\bigr]Q(\lambda)
-(-\alpha)
\det
\begin{pmatrix}
-1&-1\\
0&\lambda+\gamma
\end{pmatrix}\nonumber\\
&=\bigl[\lambda+\alpha(1+\ell)\bigr]Q(\lambda)
-\alpha(\lambda+\gamma).
\label{eq:determinante_transferencia}
\end{align}

Para calcular la primera componente de
\((\lambda I-P)^{-1}b\), la regla de Cramer sustituye la primera columna de
\cref{eq:matriz_resolvente} por
\(b=(-\alpha,0,0)^{\T}\). El determinante del numerador es
\[
\det
\begin{pmatrix}
-\alpha&-\alpha&0\\
0&\lambda+1&-1\\
0&\beta&\lambda+\gamma
\end{pmatrix}
=-\alpha Q(\lambda).
\]
Como \(r^{\T}\) selecciona esa primera componente,
\begin{equation}
W_q(s)
=r^{\T}(s^qI-P)^{-1}b
=-\alpha\frac{Q(s^q)}{\Delta_\ell(s^q)}.
\label{eq:transferencia_cramer}
\end{equation}

La identidad
\[
P-s^qI=-(s^qI-P)
\]
implica
\[
r^{\T}(P-s^qI)^{-1}b=-W_q(s).
\]
Por ello, cambiar el orden de los términos dentro de la resolvente cambia
simultáneamente el signo de la transferencia y el signo del cierre. No cambia
la solución física si ambas modificaciones se realizan de manera coherente.

\fi
\ifnum\HAFOderivationpart=3\relax
\subsection{Derivación completa de los equilibrios}
\label{app:equilibrios}

En un equilibrio \(e=(x_e,y_e,z_e)^{\T}\), las dos últimas ecuaciones de
ambos sistemas dan
\[
0=x_e-y_e+z_e,\qquad
0=-\beta y_e-\gamma z_e.
\]
De la primera,
\[
z_e=y_e-x_e.
\]
Al sustituirla en la segunda,
\begin{align*}
0
&=-\beta y_e-\gamma(y_e-x_e)\\
&=-(\beta+\gamma)y_e+\gamma x_e.
\end{align*}
Si \(\beta+\gamma\neq0\),
\begin{equation}
y_e=\frac{\gamma}{\beta+\gamma}x_e,\qquad
z_e=-\frac{\beta}{\beta+\gamma}x_e.
\label{eq:equilibrios_yz}
\end{equation}
La primera ecuación de equilibrio exige
\[
0=y_e-x_e-h(x_e).
\]
Mediante \cref{eq:equilibrios_yz},
\begin{equation}
h(x_e)+\frac{\beta}{\beta+\gamma}x_e=0.
\label{eq:equilibrio_escalar_comun}
\end{equation}
Esta única ecuación escalar determina todos los equilibrios; las otras dos
coordenadas se recuperan después con \cref{eq:equilibrios_yz}.

\subsubsection{Equilibrios del Chua no suave}

En la región \(|x_e|<1\), \(h(x_e)=m_0x_e\). Entonces
\begin{equation}
\left(m_0+\frac{\beta}{\beta+\gamma}\right)x_e=0.
\label{eq:equilibrio_interior_ns}
\end{equation}
Cuando el coeficiente no es cero, la única raíz interior es \(x_e=0\).

Para \(x_e>1\),
\[
h(x_e)=m_1x_e+(m_0-m_1),
\]
y \cref{eq:equilibrio_escalar_comun} queda
\[
\left(m_1+\frac{\beta}{\beta+\gamma}\right)x_e+(m_0-m_1)=0.
\]
Por tanto,
\begin{equation}
x_\star=
\frac{m_1-m_0}
{m_1+\dfrac{\beta}{\beta+\gamma}}.
\label{eq:raiz_exterior_ns}
\end{equation}
Si \(x_\star>1\), existe el equilibrio positivo. Como la no linealidad es
impar, existe también la raíz \(-x_\star<-1\). Los tres equilibrios son
\begin{equation}
E_0=(0,0,0),\qquad
E_\pm=
\left(
\pm x_\star,\,
\pm\frac{\gamma}{\beta+\gamma}x_\star,\,
\mp\frac{\beta}{\beta+\gamma}x_\star
\right).
\label{eq:equilibrios_ns_simbolicos}
\end{equation}
La condición regional \(x_\star>1\) debe comprobarse; no basta con resolver la
ecuación lineal exterior.

\subsubsection{Equilibrios del Chua con arctangente}

Al usar
\(h(x)=a_1x+a_2\arctan(\rho x)\) en
\cref{eq:equilibrio_escalar_comun}, se obtiene
\begin{equation}
\left(a_1+\frac{\beta}{\beta+\gamma}\right)x_e
+a_2\arctan(\rho x_e)=0.
\label{eq:equilibrio_escalar_arctan}
\end{equation}
La expresión es impar, de modo que \(x_e=0\) siempre es raíz y toda raíz
positiva no nula tiene una compañera negativa. Las raíces no nulas se
determinan numéricamente y se sustituyen en \cref{eq:equilibrios_yz}. Esta
ecuación se evalúa con \(\rho=1\) para el caso bibliográfico y con el valor de
\(\rho\) de \cref{eq:parametros_c590} para el caso \texttt{c590}; no deben
intercambiarse las dos parametrizaciones.

\fi
\ifnum\HAFOderivationpart=4\relax
\subsection{Derivación completa de los jacobianos regionales}
\label{app:jacobianos}

Para la forma de Lur'e,
\[
F(X)=PX+b\psi(r^{\T}X).
\]
La regla de la cadena produce
\begin{equation}
J_e=P+b\,\psi'(r^{\T}e)r^{\T}.
\label{eq:jacobiano_lure}
\end{equation}
Como \(r^{\T}e=x_e\), solamente cambia la entrada \((1,1)\).

En el sistema no suave,
\[
\psi'_{\mathrm{ns}}(x)=
\begin{cases}
m_0-m_1,&|x|<1,\\
0,&|x|>1.
\end{cases}
\]
Por ello,
\begin{equation}
J_{\mathrm{in}}=
\begin{pmatrix}
-\alpha(1+m_0)&\alpha&0\\
1&-1&1\\
0&-\beta&-\gamma
\end{pmatrix},
\qquad
J_{\mathrm{out}}=
\begin{pmatrix}
-\alpha(1+m_1)&\alpha&0\\
1&-1&1\\
0&-\beta&-\gamma
\end{pmatrix}.
\label{eq:jacobianos_ns}
\end{equation}
En \(x=\pm1\) la derivada clásica no existe. Los equilibrios usados en este
reporte no se encuentran sobre esas superficies, por lo que se emplea el
jacobiano de la región que contiene a cada uno.

Para la arctangente,
\[
\psi'_{\mathrm a}(x)
=\frac{a_2\rho}{1+\rho^2x^2},
\]
y
\begin{equation}
J_e=
\begin{pmatrix}
\displaystyle
-\alpha\left(1+a_1+
\frac{a_2\rho}{1+\rho^2x_e^2}\right)
&\alpha&0\\[3mm]
1&-1&1\\
0&-\beta&-\gamma
\end{pmatrix}.
\label{eq:jacobiano_arctan}
\end{equation}
Los valores propios de
\cref{eq:jacobianos_ns,eq:jacobiano_arctan} se clasifican con
\cref{eq:matignon}; la estabilidad local no sustituye el sondeo de cuenca.

\fi
\ifnum\HAFOderivationpart=5\relax
\subsection{Derivación completa del polinomio característico}
\label{app:polinomio_caracteristico}

Sea \(g=h'(x_e)\) la pendiente total de la no linealidad en un equilibrio.
El jacobiano de cualquiera de las dos familias puede escribirse como
\begin{equation}
J(g)=
\begin{pmatrix}
-\alpha(1+g)&\alpha&0\\
1&-1&1\\
0&-\beta&-\gamma
\end{pmatrix}.
\label{eq:jacobiano_pendiente_total}
\end{equation}
Por tanto,
\[
\lambda I-J(g)=
\begin{pmatrix}
\lambda+\alpha(1+g)&-\alpha&0\\
-1&\lambda+1&-1\\
0&\beta&\lambda+\gamma
\end{pmatrix}.
\]
El desarrollo por la primera fila, igual al de
\cref{eq:determinante_transferencia}, produce
\begin{equation}
\chi_g(\lambda)
=\bigl[\lambda+\alpha(1+g)\bigr]
\bigl[(\lambda+1)(\lambda+\gamma)+\beta\bigr]
-\alpha(\lambda+\gamma).
\label{eq:polinomio_caracteristico_agrupado}
\end{equation}
Como
\[
(\lambda+1)(\lambda+\gamma)+\beta
=\lambda^2+(1+\gamma)\lambda+(\beta+\gamma),
\]
la multiplicación y agrupación por potencias de \(\lambda\) da
\begin{equation}
\begin{aligned}
\chi_g(\lambda)
={}&\lambda^3+
\bigl[1+\gamma+\alpha(1+g)\bigr]\lambda^2\\
&+\bigl[
\beta+\gamma+\alpha(1+g)(1+\gamma)-\alpha
\bigr]\lambda\\
&+\alpha\bigl[(1+g)(\beta+\gamma)-\gamma\bigr].
\end{aligned}
\label{eq:polinomio_caracteristico_expandido}
\end{equation}
En el sistema no suave se sustituye \(g=m_0\) en \(E_0\) y \(g=m_1\)
en \(E_\pm\). En el sistema con arctangente,
\[
g=a_1+\frac{a_2\rho}{1+\rho^2x_e^2}.
\]
Las raíces de \cref{eq:polinomio_caracteristico_expandido} son los valores
propios utilizados en \cref{eq:matignon}; así, la clasificación local se
obtiene a partir del mismo campo que define la integración.

\fi
\ifnum\HAFOderivationpart=6\relax
\subsection{Derivación completa para la saturación centrada}
\label{app:df_saturacion_centrada}

Sea
\[
\psi_{\mathrm{ns}}(\sigma)=\Delta\operatorname{sat}(\sigma),
\qquad \Delta=m_0-m_1.
\]
Si \(0<A\leq1\), no hay saturación y
\[
\psi_{\mathrm{ns}}(A\sin\theta)=\Delta A\sin\theta.
\]
Como \(\int_0^\pi\sin^2\theta\dd\theta=\pi/2\),
\[
N_{\mathrm{ns}}(A)
=\frac{2}{\pi A}\Delta A\frac{\pi}{2}
=\Delta.
\]

Para \(A>1\), defínase
\[
\theta_0=\arcsin\left(\frac1A\right),
\qquad 0<\theta_0<\frac{\pi}{2}.
\]
En \(0\leq\theta\leq\theta_0\) y
\(\pi-\theta_0\leq\theta\leq\pi\), la entrada permanece en la parte lineal;
en el intervalo central la saturación vale \(1\). Entonces
\begin{align}
N_{\mathrm{ns}}(A)
&=\frac{2\Delta}{\pi A}
\left[
2A\int_0^{\theta_0}\sin^2\theta\dd\theta
+\int_{\theta_0}^{\pi-\theta_0}\sin\theta\dd\theta
\right]\nonumber\\
&=\frac{2\Delta}{\pi A}
\left[
A\theta_0-\frac{A}{2}\sin(2\theta_0)
+2\cos\theta_0
\right]\nonumber\\
&=\frac{2\Delta}{\pi A}
\left[A\theta_0+\cos\theta_0\right].
\label{eq:df_sat_pasos}
\end{align}
En el último paso se usó
\[
\frac{A}{2}\sin(2\theta_0)
=A\sin\theta_0\cos\theta_0
=\cos\theta_0.
\]
Finalmente,
\[
\cos\theta_0
=\sqrt{1-\frac1{A^2}}
=\frac{\sqrt{A^2-1}}{A},
\]
y \cref{eq:df_sat_pasos} se convierte en
\[
N_{\mathrm{ns}}(A)
=\Delta\frac{2}{\pi}
\left[
\arcsin\left(\frac1A\right)
+\frac{\sqrt{A^2-1}}{A^2}
\right],
\]
que coincide con \cref{eq:df_saturacion}.

\fi
\ifnum\HAFOderivationpart=7\relax
\subsection{Derivación completa de la saturación sesgada}
\label{app:df_saturacion_sesgada}

Considérese
\[
u(\theta)=c+A\cos\theta,\qquad A>0.
\]
Para escribir una fórmula que incluya los casos parcialmente y completamente
saturados, se define
\[
[v]_{-1}^{1}=\max(-1,\min(1,v))
\]
y los ángulos
\begin{equation}
\theta_+
=\arccos\left[\frac{1-c}{A}\right]_{-1}^{1},
\qquad
\theta_-
=\arccos\left[\frac{-1-c}{A}\right]_{-1}^{1}.
\label{eq:angulos_saturacion_sesgada}
\end{equation}
Como \((1-c)/A>(-1-c)/A\) y \(\arccos\) es decreciente,
\(0\leq\theta_+\leq\theta_-\leq\pi\). En el intervalo
\([0,\pi]\),
\[
\operatorname{sat}(u(\theta))=
\begin{cases}
1,&0\leq\theta\leq\theta_+,\\
c+A\cos\theta,&\theta_+\leq\theta\leq\theta_-,\\
-1,&\theta_-\leq\theta\leq\pi.
\end{cases}
\]
La simetría respecto de \(2\pi\) permite integrar solamente en
\([0,\pi]\). Para \(\psi=\Delta\operatorname{sat}\),
\begin{align}
\Psi_0(c,A)
=\frac{\Delta}{\pi}\Big[
&\theta_+
+c(\theta_--\theta_+)
+A(\sin\theta_--\sin\theta_+)\nonumber\\
&-(\pi-\theta_-)
\Big].
\label{eq:psi0_saturacion_cerrada}
\end{align}

Para la componente fundamental se usan
\[
\int\cos\theta\dd\theta=\sin\theta,\qquad
\int\cos^2\theta\dd\theta
=\frac{\theta}{2}+\frac{\sin2\theta}{4}.
\]
La sustitución por intervalos da
\begin{align}
\Psi_1(c,A)
=\frac{2\Delta}{\pi}\Bigg[
&\sin\theta_++\sin\theta_-
+c(\sin\theta_--\sin\theta_+)\nonumber\\
&+\frac{A}{2}(\theta_--\theta_+)
+\frac{A}{4}
\bigl(\sin2\theta_--\sin2\theta_+\bigr)
\Bigg].
\label{eq:psi1_saturacion_cerrada}
\end{align}
Las operaciones de recorte en
\cref{eq:angulos_saturacion_sesgada} hacen que
\cref{eq:psi0_saturacion_cerrada,eq:psi1_saturacion_cerrada} sigan siendo
válidas cuando alguna de las tres regiones tiene longitud cero.

La ecuación de media puede reducirse a una igualdad escalar. De
\cref{eq:transferencia_cramer},
\begin{align}
W_q(0)
&=-\alpha
\frac{\beta+\gamma}
{\alpha(1+\ell)(\beta+\gamma)-\alpha\gamma}\nonumber\\
&=-\frac{\beta+\gamma}
{(1+\ell)(\beta+\gamma)-\gamma}.
\label{eq:ganancia_dc}
\end{align}
Por tanto,
\begin{equation}
c=W_q(0)\Psi_0(c,A),
\label{eq:cierre_dc_escalar}
\end{equation}
y el cierre armónico es
\[
1-W_q(\ii\omega)\frac{\Psi_1(c,A)}{A}=0.
\]
Estas tres ecuaciones reales ---una para la media, y las partes real e
imaginaria del cierre--- determinan \((c,A,\omega)\).

\fi
\ifnum\HAFOderivationpart=8\relax
\subsection{Derivación completa para la arctangente}
\label{app:df_arctan}

Sea
\[
\psi_{\mathrm a}(\sigma)=a_2\arctan(\rho\sigma),
\qquad
\xi=\rho A.
\]
Se define
\[
I(\xi)=
\int_0^\pi\arctan(\xi\sin\theta)\sin\theta\dd\theta.
\]
Como \(I(0)=0\), puede calcularse \(I\) integrando su derivada:
\begin{align}
I'(\xi)
&=\int_0^\pi
\frac{\sin^2\theta}{1+\xi^2\sin^2\theta}\dd\theta\nonumber\\
&=\frac1{\xi^2}
\left[
\pi-\int_0^\pi
\frac{\dd\theta}{1+\xi^2\sin^2\theta}
\right]\nonumber\\
&=\frac{\pi}{\xi^2}
\left(1-\frac1{\sqrt{1+\xi^2}}\right).
\label{eq:derivada_integral_arctan}
\end{align}
En la última igualdad se usó la integral estándar
\[
\int_0^\pi\frac{\dd\theta}{1+\xi^2\sin^2\theta}
=\frac{\pi}{\sqrt{1+\xi^2}}.
\]
Una primitiva que se anula en \(\xi=0\) es
\begin{equation}
I(\xi)
=\pi\frac{\sqrt{1+\xi^2}-1}{\xi}.
\label{eq:integral_arctan}
\end{equation}
Al sustituir \(\xi=\rho A\) en
\cref{eq:df_centrada},
\begin{align}
N_{\mathrm a}(A)
&=\frac{2a_2}{\pi A}I(\rho A)\nonumber\\
&=\frac{2a_2}{\rho A^2}
\left(\sqrt{1+\rho^2A^2}-1\right)\nonumber\\
&=\frac{2a_2\rho}
{1+\sqrt{1+\rho^2A^2}}.
\label{eq:df_arctan_derivada}
\end{align}

Si el cierre exige \(N_{\mathrm a}(A)=k\), la amplitud puede despejarse.
De la forma racionalizada,
\[
1+\sqrt{1+\rho^2A^2}=\frac{2a_2\rho}{k}.
\]
Al aislar la raíz, elevar al cuadrado y simplificar,
\begin{equation}
A^2
=\frac{4a_2(a_2\rho-k)}{k^2\rho}.
\label{eq:amplitud_arctan}
\end{equation}
La solución se acepta solamente si el lado derecho es positivo y la igualdad
previa a elevar al cuadrado tiene lado derecho no menor que \(2\). De manera
equivalente, para \(A>0\), \(k\) tiene el mismo signo que \(a_2\) y se
encuentra estrictamente entre \(0\) y \(a_2\rho\), con el orden de los
extremos ajustado al signo de \(a_2\).

\fi
\ifnum\HAFOderivationpart=9\relax
\subsection{Derivación completa de la identidad modal y su normalización}
\label{app:identidad_modal}

Sea \(P_0=P+bkr^{\T}\). Entonces
\[
\zeta I-P_0=(\zeta I-P)-bkr^{\T}.
\]
Si \(\zeta I-P\) es invertible, el lema del determinante matricial
\(\det(M+uv^{\T})=\det(M)(1+v^{\T}M^{-1}u)\), con
\[
M=\zeta I-P,\qquad u=-bk,\qquad v=r,
\]
produce
\begin{align}
\det(\zeta I-P_0)
&=\det(\zeta I-P)
\left[1-kr^{\T}(\zeta I-P)^{-1}b\right]\nonumber\\
&=\det(\zeta I-P)\bigl[1-kG(\zeta)\bigr],
\label{eq:lema_determinante_aplicado}
\end{align}
donde
\[
G(\zeta)=r^{\T}(\zeta I-P)^{-1}b,
\qquad W_q(s)=G(s^q).
\]
Por continuidad, la identidad polinómica se extiende también a los valores en
los que la resolvente de \(P\) es singular. Si
\[
1-kG(\zeta_0)=0,
\]
entonces \(\det(\zeta_0I-P_0)=0\), de modo que existe
\(v\neq0\) con \(P_0v=\zeta_0v\). Siempre que \(r^{\T}v\neq0\), se reemplaza
\(v\) por \(v/(r^{\T}v)\) y se obtiene \(r^{\T}v=1\). Ésta es la
normalización usada en \cref{eq:modo_normalizado,eq:semilla_modal}.

En la rama sesgada,
\[
X_1=\bigl[(\ii\omega_0)^qI-P\bigr]^{-1}b\Psi_1.
\]
Al premultiplicar por \(r^{\T}\),
\[
r^{\T}X_1
=W_q(\ii\omega_0)\Psi_1
=W_q(\ii\omega_0)N_1A.
\]
El cierre \(1-W_qN_1=0\) prueba algebraicamente que
\(r^{\T}X_1=A\).

\fi
\ifnum\HAFOderivationpart=10\relax
\subsection{Derivación completa de la homotopía afín}
\label{app:homotopia_afin}

La aproximación sesgada separa
\[
\psi(\sigma)
=\Psi_0+N_1(\sigma-c)
+\bigl[\psi(\sigma)-\Psi_0-N_1(\sigma-c)\bigr].
\]
Se definen
\begin{equation}
P_{\mathrm a}=P+N_1br^{\T},\qquad
d_{\mathrm a}=b(\Psi_0-N_1c).
\label{eq:parametros_homotopia_afin}
\end{equation}
La familia de campos es
\begin{equation}
\begin{aligned}
F_\varepsilon(X)
&=P_{\mathrm a}X+d_{\mathrm a}\\
&\quad+\varepsilon b
\left[
\psi(r^{\T}X)-\Psi_0-N_1(r^{\T}X-c)
\right],
\qquad 0\leq\varepsilon\leq1.
\end{aligned}
\label{eq:homotopia_afin}
\end{equation}

En \(\varepsilon=0\), la media \(\bar X\) de
\cref{eq:cierre_dc_sesgado} es un equilibrio del auxiliar:
\begin{align*}
F_0(\bar X)
&=P\bar X+bN_1c+b\Psi_0-bN_1c\\
&=P\bar X+b\Psi_0=0.
\end{align*}
Para una perturbación \(\xi=X-\bar X\),
\[
F_0(\bar X+\xi)=P_{\mathrm a}\xi,
\]
de modo que el auxiliar conserva tanto la media como el modo armónico.

En \(\varepsilon=1\), al agrupar términos,
\begin{align*}
F_1(X)
&=PX+bN_1r^{\T}X+b\Psi_0-bN_1c\\
&\quad+b\psi(r^{\T}X)-b\Psi_0-bN_1r^{\T}X+bN_1c\\
&=PX+b\psi(r^{\T}X).
\end{align*}
Así, el extremo final de \cref{eq:homotopia_afin} es exactamente el sistema
físico de \cref{eq:lure_comun}. Una homotopía centrada de la forma
\(P_0X+\varepsilon b[\psi(\sigma)-k\sigma]\) no conserva, en general, la
media \(c\neq0\); por eso no se utiliza para la rama sesgada.
\fi

\endgroup
Por tanto,
\begin{equation}
N_{\mathrm{ns}}(A)=
\begin{cases}
\Delta,&0<A\leq1,\\[1mm]
\displaystyle
\Delta\frac{2}{\pi}
\left[
\arcsin\!\left(\frac1A\right)
+\frac{\sqrt{A^2-1}}{A^2}
\right],&A>1.
\end{cases}
\label{eq:df_saturacion}
\end{equation}

Para la arctangente, se calcula primero la integral y se justifican la
racionalización y las condiciones del despeje de amplitud:
\begingroup
\def\HAFOderivationpart{8}
% Este archivo es una biblioteca de derivaciones. Cada bloque se inserta en el
% lugar del cuerpo donde se usa, mediante \HAFOderivationpart. No se imprime
% como apéndice independiente.
\ifnum\HAFOderivationpart=1\relax
\subsection{Derivación completa de la forma de Lur'e}
\label{app:verificacion_lure}

Sea \(h(x)=\ell x+\psi(x)\). La primera ecuación de los dos modelos es
\[
{}^{\mathrm C}D_{t_0}^{q}x
=\alpha(y-x-h(x)).
\]
Al distribuir \(\alpha\),
\begin{align}
\alpha(y-x-h(x))
&=\alpha y-\alpha x-\alpha\ell x-\alpha\psi(x)\nonumber\\
&=-\alpha(1+\ell)x+\alpha y-\alpha\psi(x).
\label{eq:expansion_primera_fila}
\end{align}
Por otra parte,
\[
PX=
\begin{pmatrix}
-\alpha(1+\ell)x+\alpha y\\
x-y+z\\
-\beta y-\gamma z
\end{pmatrix},
\qquad
b\psi(r^{\T}X)=
\begin{pmatrix}
-\alpha\psi(x)\\0\\0
\end{pmatrix}.
\]
La suma reproduce \cref{eq:expansion_primera_fila} y las otras dos
ecuaciones. Para la saturación se toma
\(\ell=m_1\) y
\(\psi=(m_0-m_1)\operatorname{sat}\); para la arctangente,
\(\ell=a_1\) y \(\psi=a_2\arctan(\rho\,\cdot)\).

\fi
\ifnum\HAFOderivationpart=2\relax
\subsection{Derivación completa del determinante y la transferencia}
\label{app:transferencia_explicita}

Con \(\lambda=s^q\),
\begin{equation}
\lambda I-P=
\begin{pmatrix}
\lambda+\alpha(1+\ell)&-\alpha&0\\
-1&\lambda+1&-1\\
0&\beta&\lambda+\gamma
\end{pmatrix}.
\label{eq:matriz_resolvente}
\end{equation}
El menor de la posición \((1,1)\) es
\begin{align}
Q(\lambda)
&=
\det
\begin{pmatrix}
\lambda+1&-1\\
\beta&\lambda+\gamma
\end{pmatrix}\nonumber\\
&=(\lambda+1)(\lambda+\gamma)+\beta.
\label{eq:menor_Q}
\end{align}
Al desarrollar \(\det(\lambda I-P)\) por la primera fila,
\begin{align}
\Delta_\ell(\lambda)
&=\bigl[\lambda+\alpha(1+\ell)\bigr]Q(\lambda)
-(-\alpha)
\det
\begin{pmatrix}
-1&-1\\
0&\lambda+\gamma
\end{pmatrix}\nonumber\\
&=\bigl[\lambda+\alpha(1+\ell)\bigr]Q(\lambda)
-\alpha(\lambda+\gamma).
\label{eq:determinante_transferencia}
\end{align}

Para calcular la primera componente de
\((\lambda I-P)^{-1}b\), la regla de Cramer sustituye la primera columna de
\cref{eq:matriz_resolvente} por
\(b=(-\alpha,0,0)^{\T}\). El determinante del numerador es
\[
\det
\begin{pmatrix}
-\alpha&-\alpha&0\\
0&\lambda+1&-1\\
0&\beta&\lambda+\gamma
\end{pmatrix}
=-\alpha Q(\lambda).
\]
Como \(r^{\T}\) selecciona esa primera componente,
\begin{equation}
W_q(s)
=r^{\T}(s^qI-P)^{-1}b
=-\alpha\frac{Q(s^q)}{\Delta_\ell(s^q)}.
\label{eq:transferencia_cramer}
\end{equation}

La identidad
\[
P-s^qI=-(s^qI-P)
\]
implica
\[
r^{\T}(P-s^qI)^{-1}b=-W_q(s).
\]
Por ello, cambiar el orden de los términos dentro de la resolvente cambia
simultáneamente el signo de la transferencia y el signo del cierre. No cambia
la solución física si ambas modificaciones se realizan de manera coherente.

\fi
\ifnum\HAFOderivationpart=3\relax
\subsection{Derivación completa de los equilibrios}
\label{app:equilibrios}

En un equilibrio \(e=(x_e,y_e,z_e)^{\T}\), las dos últimas ecuaciones de
ambos sistemas dan
\[
0=x_e-y_e+z_e,\qquad
0=-\beta y_e-\gamma z_e.
\]
De la primera,
\[
z_e=y_e-x_e.
\]
Al sustituirla en la segunda,
\begin{align*}
0
&=-\beta y_e-\gamma(y_e-x_e)\\
&=-(\beta+\gamma)y_e+\gamma x_e.
\end{align*}
Si \(\beta+\gamma\neq0\),
\begin{equation}
y_e=\frac{\gamma}{\beta+\gamma}x_e,\qquad
z_e=-\frac{\beta}{\beta+\gamma}x_e.
\label{eq:equilibrios_yz}
\end{equation}
La primera ecuación de equilibrio exige
\[
0=y_e-x_e-h(x_e).
\]
Mediante \cref{eq:equilibrios_yz},
\begin{equation}
h(x_e)+\frac{\beta}{\beta+\gamma}x_e=0.
\label{eq:equilibrio_escalar_comun}
\end{equation}
Esta única ecuación escalar determina todos los equilibrios; las otras dos
coordenadas se recuperan después con \cref{eq:equilibrios_yz}.

\subsubsection{Equilibrios del Chua no suave}

En la región \(|x_e|<1\), \(h(x_e)=m_0x_e\). Entonces
\begin{equation}
\left(m_0+\frac{\beta}{\beta+\gamma}\right)x_e=0.
\label{eq:equilibrio_interior_ns}
\end{equation}
Cuando el coeficiente no es cero, la única raíz interior es \(x_e=0\).

Para \(x_e>1\),
\[
h(x_e)=m_1x_e+(m_0-m_1),
\]
y \cref{eq:equilibrio_escalar_comun} queda
\[
\left(m_1+\frac{\beta}{\beta+\gamma}\right)x_e+(m_0-m_1)=0.
\]
Por tanto,
\begin{equation}
x_\star=
\frac{m_1-m_0}
{m_1+\dfrac{\beta}{\beta+\gamma}}.
\label{eq:raiz_exterior_ns}
\end{equation}
Si \(x_\star>1\), existe el equilibrio positivo. Como la no linealidad es
impar, existe también la raíz \(-x_\star<-1\). Los tres equilibrios son
\begin{equation}
E_0=(0,0,0),\qquad
E_\pm=
\left(
\pm x_\star,\,
\pm\frac{\gamma}{\beta+\gamma}x_\star,\,
\mp\frac{\beta}{\beta+\gamma}x_\star
\right).
\label{eq:equilibrios_ns_simbolicos}
\end{equation}
La condición regional \(x_\star>1\) debe comprobarse; no basta con resolver la
ecuación lineal exterior.

\subsubsection{Equilibrios del Chua con arctangente}

Al usar
\(h(x)=a_1x+a_2\arctan(\rho x)\) en
\cref{eq:equilibrio_escalar_comun}, se obtiene
\begin{equation}
\left(a_1+\frac{\beta}{\beta+\gamma}\right)x_e
+a_2\arctan(\rho x_e)=0.
\label{eq:equilibrio_escalar_arctan}
\end{equation}
La expresión es impar, de modo que \(x_e=0\) siempre es raíz y toda raíz
positiva no nula tiene una compañera negativa. Las raíces no nulas se
determinan numéricamente y se sustituyen en \cref{eq:equilibrios_yz}. Esta
ecuación se evalúa con \(\rho=1\) para el caso bibliográfico y con el valor de
\(\rho\) de \cref{eq:parametros_c590} para el caso \texttt{c590}; no deben
intercambiarse las dos parametrizaciones.

\fi
\ifnum\HAFOderivationpart=4\relax
\subsection{Derivación completa de los jacobianos regionales}
\label{app:jacobianos}

Para la forma de Lur'e,
\[
F(X)=PX+b\psi(r^{\T}X).
\]
La regla de la cadena produce
\begin{equation}
J_e=P+b\,\psi'(r^{\T}e)r^{\T}.
\label{eq:jacobiano_lure}
\end{equation}
Como \(r^{\T}e=x_e\), solamente cambia la entrada \((1,1)\).

En el sistema no suave,
\[
\psi'_{\mathrm{ns}}(x)=
\begin{cases}
m_0-m_1,&|x|<1,\\
0,&|x|>1.
\end{cases}
\]
Por ello,
\begin{equation}
J_{\mathrm{in}}=
\begin{pmatrix}
-\alpha(1+m_0)&\alpha&0\\
1&-1&1\\
0&-\beta&-\gamma
\end{pmatrix},
\qquad
J_{\mathrm{out}}=
\begin{pmatrix}
-\alpha(1+m_1)&\alpha&0\\
1&-1&1\\
0&-\beta&-\gamma
\end{pmatrix}.
\label{eq:jacobianos_ns}
\end{equation}
En \(x=\pm1\) la derivada clásica no existe. Los equilibrios usados en este
reporte no se encuentran sobre esas superficies, por lo que se emplea el
jacobiano de la región que contiene a cada uno.

Para la arctangente,
\[
\psi'_{\mathrm a}(x)
=\frac{a_2\rho}{1+\rho^2x^2},
\]
y
\begin{equation}
J_e=
\begin{pmatrix}
\displaystyle
-\alpha\left(1+a_1+
\frac{a_2\rho}{1+\rho^2x_e^2}\right)
&\alpha&0\\[3mm]
1&-1&1\\
0&-\beta&-\gamma
\end{pmatrix}.
\label{eq:jacobiano_arctan}
\end{equation}
Los valores propios de
\cref{eq:jacobianos_ns,eq:jacobiano_arctan} se clasifican con
\cref{eq:matignon}; la estabilidad local no sustituye el sondeo de cuenca.

\fi
\ifnum\HAFOderivationpart=5\relax
\subsection{Derivación completa del polinomio característico}
\label{app:polinomio_caracteristico}

Sea \(g=h'(x_e)\) la pendiente total de la no linealidad en un equilibrio.
El jacobiano de cualquiera de las dos familias puede escribirse como
\begin{equation}
J(g)=
\begin{pmatrix}
-\alpha(1+g)&\alpha&0\\
1&-1&1\\
0&-\beta&-\gamma
\end{pmatrix}.
\label{eq:jacobiano_pendiente_total}
\end{equation}
Por tanto,
\[
\lambda I-J(g)=
\begin{pmatrix}
\lambda+\alpha(1+g)&-\alpha&0\\
-1&\lambda+1&-1\\
0&\beta&\lambda+\gamma
\end{pmatrix}.
\]
El desarrollo por la primera fila, igual al de
\cref{eq:determinante_transferencia}, produce
\begin{equation}
\chi_g(\lambda)
=\bigl[\lambda+\alpha(1+g)\bigr]
\bigl[(\lambda+1)(\lambda+\gamma)+\beta\bigr]
-\alpha(\lambda+\gamma).
\label{eq:polinomio_caracteristico_agrupado}
\end{equation}
Como
\[
(\lambda+1)(\lambda+\gamma)+\beta
=\lambda^2+(1+\gamma)\lambda+(\beta+\gamma),
\]
la multiplicación y agrupación por potencias de \(\lambda\) da
\begin{equation}
\begin{aligned}
\chi_g(\lambda)
={}&\lambda^3+
\bigl[1+\gamma+\alpha(1+g)\bigr]\lambda^2\\
&+\bigl[
\beta+\gamma+\alpha(1+g)(1+\gamma)-\alpha
\bigr]\lambda\\
&+\alpha\bigl[(1+g)(\beta+\gamma)-\gamma\bigr].
\end{aligned}
\label{eq:polinomio_caracteristico_expandido}
\end{equation}
En el sistema no suave se sustituye \(g=m_0\) en \(E_0\) y \(g=m_1\)
en \(E_\pm\). En el sistema con arctangente,
\[
g=a_1+\frac{a_2\rho}{1+\rho^2x_e^2}.
\]
Las raíces de \cref{eq:polinomio_caracteristico_expandido} son los valores
propios utilizados en \cref{eq:matignon}; así, la clasificación local se
obtiene a partir del mismo campo que define la integración.

\fi
\ifnum\HAFOderivationpart=6\relax
\subsection{Derivación completa para la saturación centrada}
\label{app:df_saturacion_centrada}

Sea
\[
\psi_{\mathrm{ns}}(\sigma)=\Delta\operatorname{sat}(\sigma),
\qquad \Delta=m_0-m_1.
\]
Si \(0<A\leq1\), no hay saturación y
\[
\psi_{\mathrm{ns}}(A\sin\theta)=\Delta A\sin\theta.
\]
Como \(\int_0^\pi\sin^2\theta\dd\theta=\pi/2\),
\[
N_{\mathrm{ns}}(A)
=\frac{2}{\pi A}\Delta A\frac{\pi}{2}
=\Delta.
\]

Para \(A>1\), defínase
\[
\theta_0=\arcsin\left(\frac1A\right),
\qquad 0<\theta_0<\frac{\pi}{2}.
\]
En \(0\leq\theta\leq\theta_0\) y
\(\pi-\theta_0\leq\theta\leq\pi\), la entrada permanece en la parte lineal;
en el intervalo central la saturación vale \(1\). Entonces
\begin{align}
N_{\mathrm{ns}}(A)
&=\frac{2\Delta}{\pi A}
\left[
2A\int_0^{\theta_0}\sin^2\theta\dd\theta
+\int_{\theta_0}^{\pi-\theta_0}\sin\theta\dd\theta
\right]\nonumber\\
&=\frac{2\Delta}{\pi A}
\left[
A\theta_0-\frac{A}{2}\sin(2\theta_0)
+2\cos\theta_0
\right]\nonumber\\
&=\frac{2\Delta}{\pi A}
\left[A\theta_0+\cos\theta_0\right].
\label{eq:df_sat_pasos}
\end{align}
En el último paso se usó
\[
\frac{A}{2}\sin(2\theta_0)
=A\sin\theta_0\cos\theta_0
=\cos\theta_0.
\]
Finalmente,
\[
\cos\theta_0
=\sqrt{1-\frac1{A^2}}
=\frac{\sqrt{A^2-1}}{A},
\]
y \cref{eq:df_sat_pasos} se convierte en
\[
N_{\mathrm{ns}}(A)
=\Delta\frac{2}{\pi}
\left[
\arcsin\left(\frac1A\right)
+\frac{\sqrt{A^2-1}}{A^2}
\right],
\]
que coincide con \cref{eq:df_saturacion}.

\fi
\ifnum\HAFOderivationpart=7\relax
\subsection{Derivación completa de la saturación sesgada}
\label{app:df_saturacion_sesgada}

Considérese
\[
u(\theta)=c+A\cos\theta,\qquad A>0.
\]
Para escribir una fórmula que incluya los casos parcialmente y completamente
saturados, se define
\[
[v]_{-1}^{1}=\max(-1,\min(1,v))
\]
y los ángulos
\begin{equation}
\theta_+
=\arccos\left[\frac{1-c}{A}\right]_{-1}^{1},
\qquad
\theta_-
=\arccos\left[\frac{-1-c}{A}\right]_{-1}^{1}.
\label{eq:angulos_saturacion_sesgada}
\end{equation}
Como \((1-c)/A>(-1-c)/A\) y \(\arccos\) es decreciente,
\(0\leq\theta_+\leq\theta_-\leq\pi\). En el intervalo
\([0,\pi]\),
\[
\operatorname{sat}(u(\theta))=
\begin{cases}
1,&0\leq\theta\leq\theta_+,\\
c+A\cos\theta,&\theta_+\leq\theta\leq\theta_-,\\
-1,&\theta_-\leq\theta\leq\pi.
\end{cases}
\]
La simetría respecto de \(2\pi\) permite integrar solamente en
\([0,\pi]\). Para \(\psi=\Delta\operatorname{sat}\),
\begin{align}
\Psi_0(c,A)
=\frac{\Delta}{\pi}\Big[
&\theta_+
+c(\theta_--\theta_+)
+A(\sin\theta_--\sin\theta_+)\nonumber\\
&-(\pi-\theta_-)
\Big].
\label{eq:psi0_saturacion_cerrada}
\end{align}

Para la componente fundamental se usan
\[
\int\cos\theta\dd\theta=\sin\theta,\qquad
\int\cos^2\theta\dd\theta
=\frac{\theta}{2}+\frac{\sin2\theta}{4}.
\]
La sustitución por intervalos da
\begin{align}
\Psi_1(c,A)
=\frac{2\Delta}{\pi}\Bigg[
&\sin\theta_++\sin\theta_-
+c(\sin\theta_--\sin\theta_+)\nonumber\\
&+\frac{A}{2}(\theta_--\theta_+)
+\frac{A}{4}
\bigl(\sin2\theta_--\sin2\theta_+\bigr)
\Bigg].
\label{eq:psi1_saturacion_cerrada}
\end{align}
Las operaciones de recorte en
\cref{eq:angulos_saturacion_sesgada} hacen que
\cref{eq:psi0_saturacion_cerrada,eq:psi1_saturacion_cerrada} sigan siendo
válidas cuando alguna de las tres regiones tiene longitud cero.

La ecuación de media puede reducirse a una igualdad escalar. De
\cref{eq:transferencia_cramer},
\begin{align}
W_q(0)
&=-\alpha
\frac{\beta+\gamma}
{\alpha(1+\ell)(\beta+\gamma)-\alpha\gamma}\nonumber\\
&=-\frac{\beta+\gamma}
{(1+\ell)(\beta+\gamma)-\gamma}.
\label{eq:ganancia_dc}
\end{align}
Por tanto,
\begin{equation}
c=W_q(0)\Psi_0(c,A),
\label{eq:cierre_dc_escalar}
\end{equation}
y el cierre armónico es
\[
1-W_q(\ii\omega)\frac{\Psi_1(c,A)}{A}=0.
\]
Estas tres ecuaciones reales ---una para la media, y las partes real e
imaginaria del cierre--- determinan \((c,A,\omega)\).

\fi
\ifnum\HAFOderivationpart=8\relax
\subsection{Derivación completa para la arctangente}
\label{app:df_arctan}

Sea
\[
\psi_{\mathrm a}(\sigma)=a_2\arctan(\rho\sigma),
\qquad
\xi=\rho A.
\]
Se define
\[
I(\xi)=
\int_0^\pi\arctan(\xi\sin\theta)\sin\theta\dd\theta.
\]
Como \(I(0)=0\), puede calcularse \(I\) integrando su derivada:
\begin{align}
I'(\xi)
&=\int_0^\pi
\frac{\sin^2\theta}{1+\xi^2\sin^2\theta}\dd\theta\nonumber\\
&=\frac1{\xi^2}
\left[
\pi-\int_0^\pi
\frac{\dd\theta}{1+\xi^2\sin^2\theta}
\right]\nonumber\\
&=\frac{\pi}{\xi^2}
\left(1-\frac1{\sqrt{1+\xi^2}}\right).
\label{eq:derivada_integral_arctan}
\end{align}
En la última igualdad se usó la integral estándar
\[
\int_0^\pi\frac{\dd\theta}{1+\xi^2\sin^2\theta}
=\frac{\pi}{\sqrt{1+\xi^2}}.
\]
Una primitiva que se anula en \(\xi=0\) es
\begin{equation}
I(\xi)
=\pi\frac{\sqrt{1+\xi^2}-1}{\xi}.
\label{eq:integral_arctan}
\end{equation}
Al sustituir \(\xi=\rho A\) en
\cref{eq:df_centrada},
\begin{align}
N_{\mathrm a}(A)
&=\frac{2a_2}{\pi A}I(\rho A)\nonumber\\
&=\frac{2a_2}{\rho A^2}
\left(\sqrt{1+\rho^2A^2}-1\right)\nonumber\\
&=\frac{2a_2\rho}
{1+\sqrt{1+\rho^2A^2}}.
\label{eq:df_arctan_derivada}
\end{align}

Si el cierre exige \(N_{\mathrm a}(A)=k\), la amplitud puede despejarse.
De la forma racionalizada,
\[
1+\sqrt{1+\rho^2A^2}=\frac{2a_2\rho}{k}.
\]
Al aislar la raíz, elevar al cuadrado y simplificar,
\begin{equation}
A^2
=\frac{4a_2(a_2\rho-k)}{k^2\rho}.
\label{eq:amplitud_arctan}
\end{equation}
La solución se acepta solamente si el lado derecho es positivo y la igualdad
previa a elevar al cuadrado tiene lado derecho no menor que \(2\). De manera
equivalente, para \(A>0\), \(k\) tiene el mismo signo que \(a_2\) y se
encuentra estrictamente entre \(0\) y \(a_2\rho\), con el orden de los
extremos ajustado al signo de \(a_2\).

\fi
\ifnum\HAFOderivationpart=9\relax
\subsection{Derivación completa de la identidad modal y su normalización}
\label{app:identidad_modal}

Sea \(P_0=P+bkr^{\T}\). Entonces
\[
\zeta I-P_0=(\zeta I-P)-bkr^{\T}.
\]
Si \(\zeta I-P\) es invertible, el lema del determinante matricial
\(\det(M+uv^{\T})=\det(M)(1+v^{\T}M^{-1}u)\), con
\[
M=\zeta I-P,\qquad u=-bk,\qquad v=r,
\]
produce
\begin{align}
\det(\zeta I-P_0)
&=\det(\zeta I-P)
\left[1-kr^{\T}(\zeta I-P)^{-1}b\right]\nonumber\\
&=\det(\zeta I-P)\bigl[1-kG(\zeta)\bigr],
\label{eq:lema_determinante_aplicado}
\end{align}
donde
\[
G(\zeta)=r^{\T}(\zeta I-P)^{-1}b,
\qquad W_q(s)=G(s^q).
\]
Por continuidad, la identidad polinómica se extiende también a los valores en
los que la resolvente de \(P\) es singular. Si
\[
1-kG(\zeta_0)=0,
\]
entonces \(\det(\zeta_0I-P_0)=0\), de modo que existe
\(v\neq0\) con \(P_0v=\zeta_0v\). Siempre que \(r^{\T}v\neq0\), se reemplaza
\(v\) por \(v/(r^{\T}v)\) y se obtiene \(r^{\T}v=1\). Ésta es la
normalización usada en \cref{eq:modo_normalizado,eq:semilla_modal}.

En la rama sesgada,
\[
X_1=\bigl[(\ii\omega_0)^qI-P\bigr]^{-1}b\Psi_1.
\]
Al premultiplicar por \(r^{\T}\),
\[
r^{\T}X_1
=W_q(\ii\omega_0)\Psi_1
=W_q(\ii\omega_0)N_1A.
\]
El cierre \(1-W_qN_1=0\) prueba algebraicamente que
\(r^{\T}X_1=A\).

\fi
\ifnum\HAFOderivationpart=10\relax
\subsection{Derivación completa de la homotopía afín}
\label{app:homotopia_afin}

La aproximación sesgada separa
\[
\psi(\sigma)
=\Psi_0+N_1(\sigma-c)
+\bigl[\psi(\sigma)-\Psi_0-N_1(\sigma-c)\bigr].
\]
Se definen
\begin{equation}
P_{\mathrm a}=P+N_1br^{\T},\qquad
d_{\mathrm a}=b(\Psi_0-N_1c).
\label{eq:parametros_homotopia_afin}
\end{equation}
La familia de campos es
\begin{equation}
\begin{aligned}
F_\varepsilon(X)
&=P_{\mathrm a}X+d_{\mathrm a}\\
&\quad+\varepsilon b
\left[
\psi(r^{\T}X)-\Psi_0-N_1(r^{\T}X-c)
\right],
\qquad 0\leq\varepsilon\leq1.
\end{aligned}
\label{eq:homotopia_afin}
\end{equation}

En \(\varepsilon=0\), la media \(\bar X\) de
\cref{eq:cierre_dc_sesgado} es un equilibrio del auxiliar:
\begin{align*}
F_0(\bar X)
&=P\bar X+bN_1c+b\Psi_0-bN_1c\\
&=P\bar X+b\Psi_0=0.
\end{align*}
Para una perturbación \(\xi=X-\bar X\),
\[
F_0(\bar X+\xi)=P_{\mathrm a}\xi,
\]
de modo que el auxiliar conserva tanto la media como el modo armónico.

En \(\varepsilon=1\), al agrupar términos,
\begin{align*}
F_1(X)
&=PX+bN_1r^{\T}X+b\Psi_0-bN_1c\\
&\quad+b\psi(r^{\T}X)-b\Psi_0-bN_1r^{\T}X+bN_1c\\
&=PX+b\psi(r^{\T}X).
\end{align*}
Así, el extremo final de \cref{eq:homotopia_afin} es exactamente el sistema
físico de \cref{eq:lure_comun}. Una homotopía centrada de la forma
\(P_0X+\varepsilon b[\psi(\sigma)-k\sigma]\) no conserva, en general, la
media \(c\neq0\); por eso no se utiliza para la rama sesgada.
\fi

\endgroup
En consecuencia, para
\(\psi_{\mathrm a}(\sigma)=a_2\arctan(\rho\sigma)\), la forma que se usará es
\begin{equation}
\begin{aligned}
N_{\mathrm a}(A)
&=\frac{2a_2}{\rho A^2}
\left(\sqrt{1+\rho^2A^2}-1\right)\\
&=\frac{2a_2\rho}
{1+\sqrt{1+\rho^2A^2}}.
\end{aligned}
\label{eq:df_arctan}
\end{equation}
La segunda forma evita cancelación por resta cuando \(\rho A\) es pequeño.

\subsection{Función descriptiva sesgada y cierre de la componente continua}

Una oscilación localizada lejos del origen requiere, en general, una componente
media no nula. Para una entrada
\[
\sigma(\theta)=c+A\cos\theta
\]
se conservan la componente continua y el primer armónico:
\begin{equation}
\begin{aligned}
\Psi_0(c,A)
&=\frac{1}{2\pi}\int_0^{2\pi}
\psi(c+A\cos\theta)\dd\theta,\\
\Psi_1(c,A)
&=\frac{1}{\pi}\int_0^{2\pi}
\psi(c+A\cos\theta)\cos\theta\dd\theta,\\
N_1(c,A)&=\frac{\Psi_1(c,A)}{A}.
\end{aligned}
\label{eq:df_sesgada}
\end{equation}
La aproximación correspondiente es
\[
\psi(c+A\cos\theta)
\simeq\Psi_0(c,A)+\Psi_1(c,A)\cos\theta.
\]
La primera igualdad que debe satisfacerse es la de la media. Si
\(\det P\neq0\),
\begin{equation}
\bar X=-P^{-1}b\,\Psi_0(c,A),\qquad
c=r^{\T}\bar X=W_q(0)\Psi_0(c,A).
\label{eq:cierre_dc_sesgado}
\end{equation}
Después se impone el cierre del primer armónico:
\begin{equation}
1-W_q(\ii\omega)N_1(c,A)=0.
\label{eq:cierre_sesgado}
\end{equation}
Por consiguiente, la búsqueda sesgada resuelve simultáneamente una ecuación de
media y una ecuación compleja de balance. Usar solamente
\cref{eq:cierre_sesgado} deja indeterminado el nivel alrededor del cual
oscila la semilla.

Para la saturación, \(\Psi_0\) y \(\Psi_1\) se obtienen separando los
intervalos en los que \(c+A\cos\theta\) está por encima de \(1\), entre
\(-1\) y \(1\), o por debajo de \(-1\):
\begingroup
\def\HAFOderivationpart{7}
% Este archivo es una biblioteca de derivaciones. Cada bloque se inserta en el
% lugar del cuerpo donde se usa, mediante \HAFOderivationpart. No se imprime
% como apéndice independiente.
\ifnum\HAFOderivationpart=1\relax
\subsection{Derivación completa de la forma de Lur'e}
\label{app:verificacion_lure}

Sea \(h(x)=\ell x+\psi(x)\). La primera ecuación de los dos modelos es
\[
{}^{\mathrm C}D_{t_0}^{q}x
=\alpha(y-x-h(x)).
\]
Al distribuir \(\alpha\),
\begin{align}
\alpha(y-x-h(x))
&=\alpha y-\alpha x-\alpha\ell x-\alpha\psi(x)\nonumber\\
&=-\alpha(1+\ell)x+\alpha y-\alpha\psi(x).
\label{eq:expansion_primera_fila}
\end{align}
Por otra parte,
\[
PX=
\begin{pmatrix}
-\alpha(1+\ell)x+\alpha y\\
x-y+z\\
-\beta y-\gamma z
\end{pmatrix},
\qquad
b\psi(r^{\T}X)=
\begin{pmatrix}
-\alpha\psi(x)\\0\\0
\end{pmatrix}.
\]
La suma reproduce \cref{eq:expansion_primera_fila} y las otras dos
ecuaciones. Para la saturación se toma
\(\ell=m_1\) y
\(\psi=(m_0-m_1)\operatorname{sat}\); para la arctangente,
\(\ell=a_1\) y \(\psi=a_2\arctan(\rho\,\cdot)\).

\fi
\ifnum\HAFOderivationpart=2\relax
\subsection{Derivación completa del determinante y la transferencia}
\label{app:transferencia_explicita}

Con \(\lambda=s^q\),
\begin{equation}
\lambda I-P=
\begin{pmatrix}
\lambda+\alpha(1+\ell)&-\alpha&0\\
-1&\lambda+1&-1\\
0&\beta&\lambda+\gamma
\end{pmatrix}.
\label{eq:matriz_resolvente}
\end{equation}
El menor de la posición \((1,1)\) es
\begin{align}
Q(\lambda)
&=
\det
\begin{pmatrix}
\lambda+1&-1\\
\beta&\lambda+\gamma
\end{pmatrix}\nonumber\\
&=(\lambda+1)(\lambda+\gamma)+\beta.
\label{eq:menor_Q}
\end{align}
Al desarrollar \(\det(\lambda I-P)\) por la primera fila,
\begin{align}
\Delta_\ell(\lambda)
&=\bigl[\lambda+\alpha(1+\ell)\bigr]Q(\lambda)
-(-\alpha)
\det
\begin{pmatrix}
-1&-1\\
0&\lambda+\gamma
\end{pmatrix}\nonumber\\
&=\bigl[\lambda+\alpha(1+\ell)\bigr]Q(\lambda)
-\alpha(\lambda+\gamma).
\label{eq:determinante_transferencia}
\end{align}

Para calcular la primera componente de
\((\lambda I-P)^{-1}b\), la regla de Cramer sustituye la primera columna de
\cref{eq:matriz_resolvente} por
\(b=(-\alpha,0,0)^{\T}\). El determinante del numerador es
\[
\det
\begin{pmatrix}
-\alpha&-\alpha&0\\
0&\lambda+1&-1\\
0&\beta&\lambda+\gamma
\end{pmatrix}
=-\alpha Q(\lambda).
\]
Como \(r^{\T}\) selecciona esa primera componente,
\begin{equation}
W_q(s)
=r^{\T}(s^qI-P)^{-1}b
=-\alpha\frac{Q(s^q)}{\Delta_\ell(s^q)}.
\label{eq:transferencia_cramer}
\end{equation}

La identidad
\[
P-s^qI=-(s^qI-P)
\]
implica
\[
r^{\T}(P-s^qI)^{-1}b=-W_q(s).
\]
Por ello, cambiar el orden de los términos dentro de la resolvente cambia
simultáneamente el signo de la transferencia y el signo del cierre. No cambia
la solución física si ambas modificaciones se realizan de manera coherente.

\fi
\ifnum\HAFOderivationpart=3\relax
\subsection{Derivación completa de los equilibrios}
\label{app:equilibrios}

En un equilibrio \(e=(x_e,y_e,z_e)^{\T}\), las dos últimas ecuaciones de
ambos sistemas dan
\[
0=x_e-y_e+z_e,\qquad
0=-\beta y_e-\gamma z_e.
\]
De la primera,
\[
z_e=y_e-x_e.
\]
Al sustituirla en la segunda,
\begin{align*}
0
&=-\beta y_e-\gamma(y_e-x_e)\\
&=-(\beta+\gamma)y_e+\gamma x_e.
\end{align*}
Si \(\beta+\gamma\neq0\),
\begin{equation}
y_e=\frac{\gamma}{\beta+\gamma}x_e,\qquad
z_e=-\frac{\beta}{\beta+\gamma}x_e.
\label{eq:equilibrios_yz}
\end{equation}
La primera ecuación de equilibrio exige
\[
0=y_e-x_e-h(x_e).
\]
Mediante \cref{eq:equilibrios_yz},
\begin{equation}
h(x_e)+\frac{\beta}{\beta+\gamma}x_e=0.
\label{eq:equilibrio_escalar_comun}
\end{equation}
Esta única ecuación escalar determina todos los equilibrios; las otras dos
coordenadas se recuperan después con \cref{eq:equilibrios_yz}.

\subsubsection{Equilibrios del Chua no suave}

En la región \(|x_e|<1\), \(h(x_e)=m_0x_e\). Entonces
\begin{equation}
\left(m_0+\frac{\beta}{\beta+\gamma}\right)x_e=0.
\label{eq:equilibrio_interior_ns}
\end{equation}
Cuando el coeficiente no es cero, la única raíz interior es \(x_e=0\).

Para \(x_e>1\),
\[
h(x_e)=m_1x_e+(m_0-m_1),
\]
y \cref{eq:equilibrio_escalar_comun} queda
\[
\left(m_1+\frac{\beta}{\beta+\gamma}\right)x_e+(m_0-m_1)=0.
\]
Por tanto,
\begin{equation}
x_\star=
\frac{m_1-m_0}
{m_1+\dfrac{\beta}{\beta+\gamma}}.
\label{eq:raiz_exterior_ns}
\end{equation}
Si \(x_\star>1\), existe el equilibrio positivo. Como la no linealidad es
impar, existe también la raíz \(-x_\star<-1\). Los tres equilibrios son
\begin{equation}
E_0=(0,0,0),\qquad
E_\pm=
\left(
\pm x_\star,\,
\pm\frac{\gamma}{\beta+\gamma}x_\star,\,
\mp\frac{\beta}{\beta+\gamma}x_\star
\right).
\label{eq:equilibrios_ns_simbolicos}
\end{equation}
La condición regional \(x_\star>1\) debe comprobarse; no basta con resolver la
ecuación lineal exterior.

\subsubsection{Equilibrios del Chua con arctangente}

Al usar
\(h(x)=a_1x+a_2\arctan(\rho x)\) en
\cref{eq:equilibrio_escalar_comun}, se obtiene
\begin{equation}
\left(a_1+\frac{\beta}{\beta+\gamma}\right)x_e
+a_2\arctan(\rho x_e)=0.
\label{eq:equilibrio_escalar_arctan}
\end{equation}
La expresión es impar, de modo que \(x_e=0\) siempre es raíz y toda raíz
positiva no nula tiene una compañera negativa. Las raíces no nulas se
determinan numéricamente y se sustituyen en \cref{eq:equilibrios_yz}. Esta
ecuación se evalúa con \(\rho=1\) para el caso bibliográfico y con el valor de
\(\rho\) de \cref{eq:parametros_c590} para el caso \texttt{c590}; no deben
intercambiarse las dos parametrizaciones.

\fi
\ifnum\HAFOderivationpart=4\relax
\subsection{Derivación completa de los jacobianos regionales}
\label{app:jacobianos}

Para la forma de Lur'e,
\[
F(X)=PX+b\psi(r^{\T}X).
\]
La regla de la cadena produce
\begin{equation}
J_e=P+b\,\psi'(r^{\T}e)r^{\T}.
\label{eq:jacobiano_lure}
\end{equation}
Como \(r^{\T}e=x_e\), solamente cambia la entrada \((1,1)\).

En el sistema no suave,
\[
\psi'_{\mathrm{ns}}(x)=
\begin{cases}
m_0-m_1,&|x|<1,\\
0,&|x|>1.
\end{cases}
\]
Por ello,
\begin{equation}
J_{\mathrm{in}}=
\begin{pmatrix}
-\alpha(1+m_0)&\alpha&0\\
1&-1&1\\
0&-\beta&-\gamma
\end{pmatrix},
\qquad
J_{\mathrm{out}}=
\begin{pmatrix}
-\alpha(1+m_1)&\alpha&0\\
1&-1&1\\
0&-\beta&-\gamma
\end{pmatrix}.
\label{eq:jacobianos_ns}
\end{equation}
En \(x=\pm1\) la derivada clásica no existe. Los equilibrios usados en este
reporte no se encuentran sobre esas superficies, por lo que se emplea el
jacobiano de la región que contiene a cada uno.

Para la arctangente,
\[
\psi'_{\mathrm a}(x)
=\frac{a_2\rho}{1+\rho^2x^2},
\]
y
\begin{equation}
J_e=
\begin{pmatrix}
\displaystyle
-\alpha\left(1+a_1+
\frac{a_2\rho}{1+\rho^2x_e^2}\right)
&\alpha&0\\[3mm]
1&-1&1\\
0&-\beta&-\gamma
\end{pmatrix}.
\label{eq:jacobiano_arctan}
\end{equation}
Los valores propios de
\cref{eq:jacobianos_ns,eq:jacobiano_arctan} se clasifican con
\cref{eq:matignon}; la estabilidad local no sustituye el sondeo de cuenca.

\fi
\ifnum\HAFOderivationpart=5\relax
\subsection{Derivación completa del polinomio característico}
\label{app:polinomio_caracteristico}

Sea \(g=h'(x_e)\) la pendiente total de la no linealidad en un equilibrio.
El jacobiano de cualquiera de las dos familias puede escribirse como
\begin{equation}
J(g)=
\begin{pmatrix}
-\alpha(1+g)&\alpha&0\\
1&-1&1\\
0&-\beta&-\gamma
\end{pmatrix}.
\label{eq:jacobiano_pendiente_total}
\end{equation}
Por tanto,
\[
\lambda I-J(g)=
\begin{pmatrix}
\lambda+\alpha(1+g)&-\alpha&0\\
-1&\lambda+1&-1\\
0&\beta&\lambda+\gamma
\end{pmatrix}.
\]
El desarrollo por la primera fila, igual al de
\cref{eq:determinante_transferencia}, produce
\begin{equation}
\chi_g(\lambda)
=\bigl[\lambda+\alpha(1+g)\bigr]
\bigl[(\lambda+1)(\lambda+\gamma)+\beta\bigr]
-\alpha(\lambda+\gamma).
\label{eq:polinomio_caracteristico_agrupado}
\end{equation}
Como
\[
(\lambda+1)(\lambda+\gamma)+\beta
=\lambda^2+(1+\gamma)\lambda+(\beta+\gamma),
\]
la multiplicación y agrupación por potencias de \(\lambda\) da
\begin{equation}
\begin{aligned}
\chi_g(\lambda)
={}&\lambda^3+
\bigl[1+\gamma+\alpha(1+g)\bigr]\lambda^2\\
&+\bigl[
\beta+\gamma+\alpha(1+g)(1+\gamma)-\alpha
\bigr]\lambda\\
&+\alpha\bigl[(1+g)(\beta+\gamma)-\gamma\bigr].
\end{aligned}
\label{eq:polinomio_caracteristico_expandido}
\end{equation}
En el sistema no suave se sustituye \(g=m_0\) en \(E_0\) y \(g=m_1\)
en \(E_\pm\). En el sistema con arctangente,
\[
g=a_1+\frac{a_2\rho}{1+\rho^2x_e^2}.
\]
Las raíces de \cref{eq:polinomio_caracteristico_expandido} son los valores
propios utilizados en \cref{eq:matignon}; así, la clasificación local se
obtiene a partir del mismo campo que define la integración.

\fi
\ifnum\HAFOderivationpart=6\relax
\subsection{Derivación completa para la saturación centrada}
\label{app:df_saturacion_centrada}

Sea
\[
\psi_{\mathrm{ns}}(\sigma)=\Delta\operatorname{sat}(\sigma),
\qquad \Delta=m_0-m_1.
\]
Si \(0<A\leq1\), no hay saturación y
\[
\psi_{\mathrm{ns}}(A\sin\theta)=\Delta A\sin\theta.
\]
Como \(\int_0^\pi\sin^2\theta\dd\theta=\pi/2\),
\[
N_{\mathrm{ns}}(A)
=\frac{2}{\pi A}\Delta A\frac{\pi}{2}
=\Delta.
\]

Para \(A>1\), defínase
\[
\theta_0=\arcsin\left(\frac1A\right),
\qquad 0<\theta_0<\frac{\pi}{2}.
\]
En \(0\leq\theta\leq\theta_0\) y
\(\pi-\theta_0\leq\theta\leq\pi\), la entrada permanece en la parte lineal;
en el intervalo central la saturación vale \(1\). Entonces
\begin{align}
N_{\mathrm{ns}}(A)
&=\frac{2\Delta}{\pi A}
\left[
2A\int_0^{\theta_0}\sin^2\theta\dd\theta
+\int_{\theta_0}^{\pi-\theta_0}\sin\theta\dd\theta
\right]\nonumber\\
&=\frac{2\Delta}{\pi A}
\left[
A\theta_0-\frac{A}{2}\sin(2\theta_0)
+2\cos\theta_0
\right]\nonumber\\
&=\frac{2\Delta}{\pi A}
\left[A\theta_0+\cos\theta_0\right].
\label{eq:df_sat_pasos}
\end{align}
En el último paso se usó
\[
\frac{A}{2}\sin(2\theta_0)
=A\sin\theta_0\cos\theta_0
=\cos\theta_0.
\]
Finalmente,
\[
\cos\theta_0
=\sqrt{1-\frac1{A^2}}
=\frac{\sqrt{A^2-1}}{A},
\]
y \cref{eq:df_sat_pasos} se convierte en
\[
N_{\mathrm{ns}}(A)
=\Delta\frac{2}{\pi}
\left[
\arcsin\left(\frac1A\right)
+\frac{\sqrt{A^2-1}}{A^2}
\right],
\]
que coincide con \cref{eq:df_saturacion}.

\fi
\ifnum\HAFOderivationpart=7\relax
\subsection{Derivación completa de la saturación sesgada}
\label{app:df_saturacion_sesgada}

Considérese
\[
u(\theta)=c+A\cos\theta,\qquad A>0.
\]
Para escribir una fórmula que incluya los casos parcialmente y completamente
saturados, se define
\[
[v]_{-1}^{1}=\max(-1,\min(1,v))
\]
y los ángulos
\begin{equation}
\theta_+
=\arccos\left[\frac{1-c}{A}\right]_{-1}^{1},
\qquad
\theta_-
=\arccos\left[\frac{-1-c}{A}\right]_{-1}^{1}.
\label{eq:angulos_saturacion_sesgada}
\end{equation}
Como \((1-c)/A>(-1-c)/A\) y \(\arccos\) es decreciente,
\(0\leq\theta_+\leq\theta_-\leq\pi\). En el intervalo
\([0,\pi]\),
\[
\operatorname{sat}(u(\theta))=
\begin{cases}
1,&0\leq\theta\leq\theta_+,\\
c+A\cos\theta,&\theta_+\leq\theta\leq\theta_-,\\
-1,&\theta_-\leq\theta\leq\pi.
\end{cases}
\]
La simetría respecto de \(2\pi\) permite integrar solamente en
\([0,\pi]\). Para \(\psi=\Delta\operatorname{sat}\),
\begin{align}
\Psi_0(c,A)
=\frac{\Delta}{\pi}\Big[
&\theta_+
+c(\theta_--\theta_+)
+A(\sin\theta_--\sin\theta_+)\nonumber\\
&-(\pi-\theta_-)
\Big].
\label{eq:psi0_saturacion_cerrada}
\end{align}

Para la componente fundamental se usan
\[
\int\cos\theta\dd\theta=\sin\theta,\qquad
\int\cos^2\theta\dd\theta
=\frac{\theta}{2}+\frac{\sin2\theta}{4}.
\]
La sustitución por intervalos da
\begin{align}
\Psi_1(c,A)
=\frac{2\Delta}{\pi}\Bigg[
&\sin\theta_++\sin\theta_-
+c(\sin\theta_--\sin\theta_+)\nonumber\\
&+\frac{A}{2}(\theta_--\theta_+)
+\frac{A}{4}
\bigl(\sin2\theta_--\sin2\theta_+\bigr)
\Bigg].
\label{eq:psi1_saturacion_cerrada}
\end{align}
Las operaciones de recorte en
\cref{eq:angulos_saturacion_sesgada} hacen que
\cref{eq:psi0_saturacion_cerrada,eq:psi1_saturacion_cerrada} sigan siendo
válidas cuando alguna de las tres regiones tiene longitud cero.

La ecuación de media puede reducirse a una igualdad escalar. De
\cref{eq:transferencia_cramer},
\begin{align}
W_q(0)
&=-\alpha
\frac{\beta+\gamma}
{\alpha(1+\ell)(\beta+\gamma)-\alpha\gamma}\nonumber\\
&=-\frac{\beta+\gamma}
{(1+\ell)(\beta+\gamma)-\gamma}.
\label{eq:ganancia_dc}
\end{align}
Por tanto,
\begin{equation}
c=W_q(0)\Psi_0(c,A),
\label{eq:cierre_dc_escalar}
\end{equation}
y el cierre armónico es
\[
1-W_q(\ii\omega)\frac{\Psi_1(c,A)}{A}=0.
\]
Estas tres ecuaciones reales ---una para la media, y las partes real e
imaginaria del cierre--- determinan \((c,A,\omega)\).

\fi
\ifnum\HAFOderivationpart=8\relax
\subsection{Derivación completa para la arctangente}
\label{app:df_arctan}

Sea
\[
\psi_{\mathrm a}(\sigma)=a_2\arctan(\rho\sigma),
\qquad
\xi=\rho A.
\]
Se define
\[
I(\xi)=
\int_0^\pi\arctan(\xi\sin\theta)\sin\theta\dd\theta.
\]
Como \(I(0)=0\), puede calcularse \(I\) integrando su derivada:
\begin{align}
I'(\xi)
&=\int_0^\pi
\frac{\sin^2\theta}{1+\xi^2\sin^2\theta}\dd\theta\nonumber\\
&=\frac1{\xi^2}
\left[
\pi-\int_0^\pi
\frac{\dd\theta}{1+\xi^2\sin^2\theta}
\right]\nonumber\\
&=\frac{\pi}{\xi^2}
\left(1-\frac1{\sqrt{1+\xi^2}}\right).
\label{eq:derivada_integral_arctan}
\end{align}
En la última igualdad se usó la integral estándar
\[
\int_0^\pi\frac{\dd\theta}{1+\xi^2\sin^2\theta}
=\frac{\pi}{\sqrt{1+\xi^2}}.
\]
Una primitiva que se anula en \(\xi=0\) es
\begin{equation}
I(\xi)
=\pi\frac{\sqrt{1+\xi^2}-1}{\xi}.
\label{eq:integral_arctan}
\end{equation}
Al sustituir \(\xi=\rho A\) en
\cref{eq:df_centrada},
\begin{align}
N_{\mathrm a}(A)
&=\frac{2a_2}{\pi A}I(\rho A)\nonumber\\
&=\frac{2a_2}{\rho A^2}
\left(\sqrt{1+\rho^2A^2}-1\right)\nonumber\\
&=\frac{2a_2\rho}
{1+\sqrt{1+\rho^2A^2}}.
\label{eq:df_arctan_derivada}
\end{align}

Si el cierre exige \(N_{\mathrm a}(A)=k\), la amplitud puede despejarse.
De la forma racionalizada,
\[
1+\sqrt{1+\rho^2A^2}=\frac{2a_2\rho}{k}.
\]
Al aislar la raíz, elevar al cuadrado y simplificar,
\begin{equation}
A^2
=\frac{4a_2(a_2\rho-k)}{k^2\rho}.
\label{eq:amplitud_arctan}
\end{equation}
La solución se acepta solamente si el lado derecho es positivo y la igualdad
previa a elevar al cuadrado tiene lado derecho no menor que \(2\). De manera
equivalente, para \(A>0\), \(k\) tiene el mismo signo que \(a_2\) y se
encuentra estrictamente entre \(0\) y \(a_2\rho\), con el orden de los
extremos ajustado al signo de \(a_2\).

\fi
\ifnum\HAFOderivationpart=9\relax
\subsection{Derivación completa de la identidad modal y su normalización}
\label{app:identidad_modal}

Sea \(P_0=P+bkr^{\T}\). Entonces
\[
\zeta I-P_0=(\zeta I-P)-bkr^{\T}.
\]
Si \(\zeta I-P\) es invertible, el lema del determinante matricial
\(\det(M+uv^{\T})=\det(M)(1+v^{\T}M^{-1}u)\), con
\[
M=\zeta I-P,\qquad u=-bk,\qquad v=r,
\]
produce
\begin{align}
\det(\zeta I-P_0)
&=\det(\zeta I-P)
\left[1-kr^{\T}(\zeta I-P)^{-1}b\right]\nonumber\\
&=\det(\zeta I-P)\bigl[1-kG(\zeta)\bigr],
\label{eq:lema_determinante_aplicado}
\end{align}
donde
\[
G(\zeta)=r^{\T}(\zeta I-P)^{-1}b,
\qquad W_q(s)=G(s^q).
\]
Por continuidad, la identidad polinómica se extiende también a los valores en
los que la resolvente de \(P\) es singular. Si
\[
1-kG(\zeta_0)=0,
\]
entonces \(\det(\zeta_0I-P_0)=0\), de modo que existe
\(v\neq0\) con \(P_0v=\zeta_0v\). Siempre que \(r^{\T}v\neq0\), se reemplaza
\(v\) por \(v/(r^{\T}v)\) y se obtiene \(r^{\T}v=1\). Ésta es la
normalización usada en \cref{eq:modo_normalizado,eq:semilla_modal}.

En la rama sesgada,
\[
X_1=\bigl[(\ii\omega_0)^qI-P\bigr]^{-1}b\Psi_1.
\]
Al premultiplicar por \(r^{\T}\),
\[
r^{\T}X_1
=W_q(\ii\omega_0)\Psi_1
=W_q(\ii\omega_0)N_1A.
\]
El cierre \(1-W_qN_1=0\) prueba algebraicamente que
\(r^{\T}X_1=A\).

\fi
\ifnum\HAFOderivationpart=10\relax
\subsection{Derivación completa de la homotopía afín}
\label{app:homotopia_afin}

La aproximación sesgada separa
\[
\psi(\sigma)
=\Psi_0+N_1(\sigma-c)
+\bigl[\psi(\sigma)-\Psi_0-N_1(\sigma-c)\bigr].
\]
Se definen
\begin{equation}
P_{\mathrm a}=P+N_1br^{\T},\qquad
d_{\mathrm a}=b(\Psi_0-N_1c).
\label{eq:parametros_homotopia_afin}
\end{equation}
La familia de campos es
\begin{equation}
\begin{aligned}
F_\varepsilon(X)
&=P_{\mathrm a}X+d_{\mathrm a}\\
&\quad+\varepsilon b
\left[
\psi(r^{\T}X)-\Psi_0-N_1(r^{\T}X-c)
\right],
\qquad 0\leq\varepsilon\leq1.
\end{aligned}
\label{eq:homotopia_afin}
\end{equation}

En \(\varepsilon=0\), la media \(\bar X\) de
\cref{eq:cierre_dc_sesgado} es un equilibrio del auxiliar:
\begin{align*}
F_0(\bar X)
&=P\bar X+bN_1c+b\Psi_0-bN_1c\\
&=P\bar X+b\Psi_0=0.
\end{align*}
Para una perturbación \(\xi=X-\bar X\),
\[
F_0(\bar X+\xi)=P_{\mathrm a}\xi,
\]
de modo que el auxiliar conserva tanto la media como el modo armónico.

En \(\varepsilon=1\), al agrupar términos,
\begin{align*}
F_1(X)
&=PX+bN_1r^{\T}X+b\Psi_0-bN_1c\\
&\quad+b\psi(r^{\T}X)-b\Psi_0-bN_1r^{\T}X+bN_1c\\
&=PX+b\psi(r^{\T}X).
\end{align*}
Así, el extremo final de \cref{eq:homotopia_afin} es exactamente el sistema
físico de \cref{eq:lure_comun}. Una homotopía centrada de la forma
\(P_0X+\varepsilon b[\psi(\sigma)-k\sigma]\) no conserva, en general, la
media \(c\neq0\); por eso no se utiliza para la rama sesgada.
\fi

\endgroup
Esta separación también permite verificar una cuadratura numérica sin evaluar
la saturación en una malla excesivamente fina.

\subsection{Construcción modal de las semillas}

Sea \(k=N_\psi(A_0)\) para una rama centrada, o
\(k=N_1(c,A)\) para una rama sesgada, y defínase
\begin{equation}
P_0=P+bkr^{\T},\qquad
\zeta_0=(\ii\omega_0)^q,\qquad
G(\zeta)=r^{\T}(\zeta I-P)^{-1}b.
\label{eq:matriz_auxiliar}
\end{equation}
La transferencia de \cref{eq:transferencia_fraccionaria} satisface
\(W_q(s)=G(s^q)\); \(s\) es la variable de Laplace y
\(\zeta=s^q\) es la variable espectral.
El lema del determinante da
\begin{equation}
\det(\zeta I-P_0)
=\det(\zeta I-P)\bigl(1-kG(\zeta)\bigr).
\label{eq:identidad_modal}
\end{equation}
La identidad y la normalización se justifican mediante el lema del
determinante matricial:
\begingroup
\def\HAFOderivationpart{9}
% Este archivo es una biblioteca de derivaciones. Cada bloque se inserta en el
% lugar del cuerpo donde se usa, mediante \HAFOderivationpart. No se imprime
% como apéndice independiente.
\ifnum\HAFOderivationpart=1\relax
\subsection{Derivación completa de la forma de Lur'e}
\label{app:verificacion_lure}

Sea \(h(x)=\ell x+\psi(x)\). La primera ecuación de los dos modelos es
\[
{}^{\mathrm C}D_{t_0}^{q}x
=\alpha(y-x-h(x)).
\]
Al distribuir \(\alpha\),
\begin{align}
\alpha(y-x-h(x))
&=\alpha y-\alpha x-\alpha\ell x-\alpha\psi(x)\nonumber\\
&=-\alpha(1+\ell)x+\alpha y-\alpha\psi(x).
\label{eq:expansion_primera_fila}
\end{align}
Por otra parte,
\[
PX=
\begin{pmatrix}
-\alpha(1+\ell)x+\alpha y\\
x-y+z\\
-\beta y-\gamma z
\end{pmatrix},
\qquad
b\psi(r^{\T}X)=
\begin{pmatrix}
-\alpha\psi(x)\\0\\0
\end{pmatrix}.
\]
La suma reproduce \cref{eq:expansion_primera_fila} y las otras dos
ecuaciones. Para la saturación se toma
\(\ell=m_1\) y
\(\psi=(m_0-m_1)\operatorname{sat}\); para la arctangente,
\(\ell=a_1\) y \(\psi=a_2\arctan(\rho\,\cdot)\).

\fi
\ifnum\HAFOderivationpart=2\relax
\subsection{Derivación completa del determinante y la transferencia}
\label{app:transferencia_explicita}

Con \(\lambda=s^q\),
\begin{equation}
\lambda I-P=
\begin{pmatrix}
\lambda+\alpha(1+\ell)&-\alpha&0\\
-1&\lambda+1&-1\\
0&\beta&\lambda+\gamma
\end{pmatrix}.
\label{eq:matriz_resolvente}
\end{equation}
El menor de la posición \((1,1)\) es
\begin{align}
Q(\lambda)
&=
\det
\begin{pmatrix}
\lambda+1&-1\\
\beta&\lambda+\gamma
\end{pmatrix}\nonumber\\
&=(\lambda+1)(\lambda+\gamma)+\beta.
\label{eq:menor_Q}
\end{align}
Al desarrollar \(\det(\lambda I-P)\) por la primera fila,
\begin{align}
\Delta_\ell(\lambda)
&=\bigl[\lambda+\alpha(1+\ell)\bigr]Q(\lambda)
-(-\alpha)
\det
\begin{pmatrix}
-1&-1\\
0&\lambda+\gamma
\end{pmatrix}\nonumber\\
&=\bigl[\lambda+\alpha(1+\ell)\bigr]Q(\lambda)
-\alpha(\lambda+\gamma).
\label{eq:determinante_transferencia}
\end{align}

Para calcular la primera componente de
\((\lambda I-P)^{-1}b\), la regla de Cramer sustituye la primera columna de
\cref{eq:matriz_resolvente} por
\(b=(-\alpha,0,0)^{\T}\). El determinante del numerador es
\[
\det
\begin{pmatrix}
-\alpha&-\alpha&0\\
0&\lambda+1&-1\\
0&\beta&\lambda+\gamma
\end{pmatrix}
=-\alpha Q(\lambda).
\]
Como \(r^{\T}\) selecciona esa primera componente,
\begin{equation}
W_q(s)
=r^{\T}(s^qI-P)^{-1}b
=-\alpha\frac{Q(s^q)}{\Delta_\ell(s^q)}.
\label{eq:transferencia_cramer}
\end{equation}

La identidad
\[
P-s^qI=-(s^qI-P)
\]
implica
\[
r^{\T}(P-s^qI)^{-1}b=-W_q(s).
\]
Por ello, cambiar el orden de los términos dentro de la resolvente cambia
simultáneamente el signo de la transferencia y el signo del cierre. No cambia
la solución física si ambas modificaciones se realizan de manera coherente.

\fi
\ifnum\HAFOderivationpart=3\relax
\subsection{Derivación completa de los equilibrios}
\label{app:equilibrios}

En un equilibrio \(e=(x_e,y_e,z_e)^{\T}\), las dos últimas ecuaciones de
ambos sistemas dan
\[
0=x_e-y_e+z_e,\qquad
0=-\beta y_e-\gamma z_e.
\]
De la primera,
\[
z_e=y_e-x_e.
\]
Al sustituirla en la segunda,
\begin{align*}
0
&=-\beta y_e-\gamma(y_e-x_e)\\
&=-(\beta+\gamma)y_e+\gamma x_e.
\end{align*}
Si \(\beta+\gamma\neq0\),
\begin{equation}
y_e=\frac{\gamma}{\beta+\gamma}x_e,\qquad
z_e=-\frac{\beta}{\beta+\gamma}x_e.
\label{eq:equilibrios_yz}
\end{equation}
La primera ecuación de equilibrio exige
\[
0=y_e-x_e-h(x_e).
\]
Mediante \cref{eq:equilibrios_yz},
\begin{equation}
h(x_e)+\frac{\beta}{\beta+\gamma}x_e=0.
\label{eq:equilibrio_escalar_comun}
\end{equation}
Esta única ecuación escalar determina todos los equilibrios; las otras dos
coordenadas se recuperan después con \cref{eq:equilibrios_yz}.

\subsubsection{Equilibrios del Chua no suave}

En la región \(|x_e|<1\), \(h(x_e)=m_0x_e\). Entonces
\begin{equation}
\left(m_0+\frac{\beta}{\beta+\gamma}\right)x_e=0.
\label{eq:equilibrio_interior_ns}
\end{equation}
Cuando el coeficiente no es cero, la única raíz interior es \(x_e=0\).

Para \(x_e>1\),
\[
h(x_e)=m_1x_e+(m_0-m_1),
\]
y \cref{eq:equilibrio_escalar_comun} queda
\[
\left(m_1+\frac{\beta}{\beta+\gamma}\right)x_e+(m_0-m_1)=0.
\]
Por tanto,
\begin{equation}
x_\star=
\frac{m_1-m_0}
{m_1+\dfrac{\beta}{\beta+\gamma}}.
\label{eq:raiz_exterior_ns}
\end{equation}
Si \(x_\star>1\), existe el equilibrio positivo. Como la no linealidad es
impar, existe también la raíz \(-x_\star<-1\). Los tres equilibrios son
\begin{equation}
E_0=(0,0,0),\qquad
E_\pm=
\left(
\pm x_\star,\,
\pm\frac{\gamma}{\beta+\gamma}x_\star,\,
\mp\frac{\beta}{\beta+\gamma}x_\star
\right).
\label{eq:equilibrios_ns_simbolicos}
\end{equation}
La condición regional \(x_\star>1\) debe comprobarse; no basta con resolver la
ecuación lineal exterior.

\subsubsection{Equilibrios del Chua con arctangente}

Al usar
\(h(x)=a_1x+a_2\arctan(\rho x)\) en
\cref{eq:equilibrio_escalar_comun}, se obtiene
\begin{equation}
\left(a_1+\frac{\beta}{\beta+\gamma}\right)x_e
+a_2\arctan(\rho x_e)=0.
\label{eq:equilibrio_escalar_arctan}
\end{equation}
La expresión es impar, de modo que \(x_e=0\) siempre es raíz y toda raíz
positiva no nula tiene una compañera negativa. Las raíces no nulas se
determinan numéricamente y se sustituyen en \cref{eq:equilibrios_yz}. Esta
ecuación se evalúa con \(\rho=1\) para el caso bibliográfico y con el valor de
\(\rho\) de \cref{eq:parametros_c590} para el caso \texttt{c590}; no deben
intercambiarse las dos parametrizaciones.

\fi
\ifnum\HAFOderivationpart=4\relax
\subsection{Derivación completa de los jacobianos regionales}
\label{app:jacobianos}

Para la forma de Lur'e,
\[
F(X)=PX+b\psi(r^{\T}X).
\]
La regla de la cadena produce
\begin{equation}
J_e=P+b\,\psi'(r^{\T}e)r^{\T}.
\label{eq:jacobiano_lure}
\end{equation}
Como \(r^{\T}e=x_e\), solamente cambia la entrada \((1,1)\).

En el sistema no suave,
\[
\psi'_{\mathrm{ns}}(x)=
\begin{cases}
m_0-m_1,&|x|<1,\\
0,&|x|>1.
\end{cases}
\]
Por ello,
\begin{equation}
J_{\mathrm{in}}=
\begin{pmatrix}
-\alpha(1+m_0)&\alpha&0\\
1&-1&1\\
0&-\beta&-\gamma
\end{pmatrix},
\qquad
J_{\mathrm{out}}=
\begin{pmatrix}
-\alpha(1+m_1)&\alpha&0\\
1&-1&1\\
0&-\beta&-\gamma
\end{pmatrix}.
\label{eq:jacobianos_ns}
\end{equation}
En \(x=\pm1\) la derivada clásica no existe. Los equilibrios usados en este
reporte no se encuentran sobre esas superficies, por lo que se emplea el
jacobiano de la región que contiene a cada uno.

Para la arctangente,
\[
\psi'_{\mathrm a}(x)
=\frac{a_2\rho}{1+\rho^2x^2},
\]
y
\begin{equation}
J_e=
\begin{pmatrix}
\displaystyle
-\alpha\left(1+a_1+
\frac{a_2\rho}{1+\rho^2x_e^2}\right)
&\alpha&0\\[3mm]
1&-1&1\\
0&-\beta&-\gamma
\end{pmatrix}.
\label{eq:jacobiano_arctan}
\end{equation}
Los valores propios de
\cref{eq:jacobianos_ns,eq:jacobiano_arctan} se clasifican con
\cref{eq:matignon}; la estabilidad local no sustituye el sondeo de cuenca.

\fi
\ifnum\HAFOderivationpart=5\relax
\subsection{Derivación completa del polinomio característico}
\label{app:polinomio_caracteristico}

Sea \(g=h'(x_e)\) la pendiente total de la no linealidad en un equilibrio.
El jacobiano de cualquiera de las dos familias puede escribirse como
\begin{equation}
J(g)=
\begin{pmatrix}
-\alpha(1+g)&\alpha&0\\
1&-1&1\\
0&-\beta&-\gamma
\end{pmatrix}.
\label{eq:jacobiano_pendiente_total}
\end{equation}
Por tanto,
\[
\lambda I-J(g)=
\begin{pmatrix}
\lambda+\alpha(1+g)&-\alpha&0\\
-1&\lambda+1&-1\\
0&\beta&\lambda+\gamma
\end{pmatrix}.
\]
El desarrollo por la primera fila, igual al de
\cref{eq:determinante_transferencia}, produce
\begin{equation}
\chi_g(\lambda)
=\bigl[\lambda+\alpha(1+g)\bigr]
\bigl[(\lambda+1)(\lambda+\gamma)+\beta\bigr]
-\alpha(\lambda+\gamma).
\label{eq:polinomio_caracteristico_agrupado}
\end{equation}
Como
\[
(\lambda+1)(\lambda+\gamma)+\beta
=\lambda^2+(1+\gamma)\lambda+(\beta+\gamma),
\]
la multiplicación y agrupación por potencias de \(\lambda\) da
\begin{equation}
\begin{aligned}
\chi_g(\lambda)
={}&\lambda^3+
\bigl[1+\gamma+\alpha(1+g)\bigr]\lambda^2\\
&+\bigl[
\beta+\gamma+\alpha(1+g)(1+\gamma)-\alpha
\bigr]\lambda\\
&+\alpha\bigl[(1+g)(\beta+\gamma)-\gamma\bigr].
\end{aligned}
\label{eq:polinomio_caracteristico_expandido}
\end{equation}
En el sistema no suave se sustituye \(g=m_0\) en \(E_0\) y \(g=m_1\)
en \(E_\pm\). En el sistema con arctangente,
\[
g=a_1+\frac{a_2\rho}{1+\rho^2x_e^2}.
\]
Las raíces de \cref{eq:polinomio_caracteristico_expandido} son los valores
propios utilizados en \cref{eq:matignon}; así, la clasificación local se
obtiene a partir del mismo campo que define la integración.

\fi
\ifnum\HAFOderivationpart=6\relax
\subsection{Derivación completa para la saturación centrada}
\label{app:df_saturacion_centrada}

Sea
\[
\psi_{\mathrm{ns}}(\sigma)=\Delta\operatorname{sat}(\sigma),
\qquad \Delta=m_0-m_1.
\]
Si \(0<A\leq1\), no hay saturación y
\[
\psi_{\mathrm{ns}}(A\sin\theta)=\Delta A\sin\theta.
\]
Como \(\int_0^\pi\sin^2\theta\dd\theta=\pi/2\),
\[
N_{\mathrm{ns}}(A)
=\frac{2}{\pi A}\Delta A\frac{\pi}{2}
=\Delta.
\]

Para \(A>1\), defínase
\[
\theta_0=\arcsin\left(\frac1A\right),
\qquad 0<\theta_0<\frac{\pi}{2}.
\]
En \(0\leq\theta\leq\theta_0\) y
\(\pi-\theta_0\leq\theta\leq\pi\), la entrada permanece en la parte lineal;
en el intervalo central la saturación vale \(1\). Entonces
\begin{align}
N_{\mathrm{ns}}(A)
&=\frac{2\Delta}{\pi A}
\left[
2A\int_0^{\theta_0}\sin^2\theta\dd\theta
+\int_{\theta_0}^{\pi-\theta_0}\sin\theta\dd\theta
\right]\nonumber\\
&=\frac{2\Delta}{\pi A}
\left[
A\theta_0-\frac{A}{2}\sin(2\theta_0)
+2\cos\theta_0
\right]\nonumber\\
&=\frac{2\Delta}{\pi A}
\left[A\theta_0+\cos\theta_0\right].
\label{eq:df_sat_pasos}
\end{align}
En el último paso se usó
\[
\frac{A}{2}\sin(2\theta_0)
=A\sin\theta_0\cos\theta_0
=\cos\theta_0.
\]
Finalmente,
\[
\cos\theta_0
=\sqrt{1-\frac1{A^2}}
=\frac{\sqrt{A^2-1}}{A},
\]
y \cref{eq:df_sat_pasos} se convierte en
\[
N_{\mathrm{ns}}(A)
=\Delta\frac{2}{\pi}
\left[
\arcsin\left(\frac1A\right)
+\frac{\sqrt{A^2-1}}{A^2}
\right],
\]
que coincide con \cref{eq:df_saturacion}.

\fi
\ifnum\HAFOderivationpart=7\relax
\subsection{Derivación completa de la saturación sesgada}
\label{app:df_saturacion_sesgada}

Considérese
\[
u(\theta)=c+A\cos\theta,\qquad A>0.
\]
Para escribir una fórmula que incluya los casos parcialmente y completamente
saturados, se define
\[
[v]_{-1}^{1}=\max(-1,\min(1,v))
\]
y los ángulos
\begin{equation}
\theta_+
=\arccos\left[\frac{1-c}{A}\right]_{-1}^{1},
\qquad
\theta_-
=\arccos\left[\frac{-1-c}{A}\right]_{-1}^{1}.
\label{eq:angulos_saturacion_sesgada}
\end{equation}
Como \((1-c)/A>(-1-c)/A\) y \(\arccos\) es decreciente,
\(0\leq\theta_+\leq\theta_-\leq\pi\). En el intervalo
\([0,\pi]\),
\[
\operatorname{sat}(u(\theta))=
\begin{cases}
1,&0\leq\theta\leq\theta_+,\\
c+A\cos\theta,&\theta_+\leq\theta\leq\theta_-,\\
-1,&\theta_-\leq\theta\leq\pi.
\end{cases}
\]
La simetría respecto de \(2\pi\) permite integrar solamente en
\([0,\pi]\). Para \(\psi=\Delta\operatorname{sat}\),
\begin{align}
\Psi_0(c,A)
=\frac{\Delta}{\pi}\Big[
&\theta_+
+c(\theta_--\theta_+)
+A(\sin\theta_--\sin\theta_+)\nonumber\\
&-(\pi-\theta_-)
\Big].
\label{eq:psi0_saturacion_cerrada}
\end{align}

Para la componente fundamental se usan
\[
\int\cos\theta\dd\theta=\sin\theta,\qquad
\int\cos^2\theta\dd\theta
=\frac{\theta}{2}+\frac{\sin2\theta}{4}.
\]
La sustitución por intervalos da
\begin{align}
\Psi_1(c,A)
=\frac{2\Delta}{\pi}\Bigg[
&\sin\theta_++\sin\theta_-
+c(\sin\theta_--\sin\theta_+)\nonumber\\
&+\frac{A}{2}(\theta_--\theta_+)
+\frac{A}{4}
\bigl(\sin2\theta_--\sin2\theta_+\bigr)
\Bigg].
\label{eq:psi1_saturacion_cerrada}
\end{align}
Las operaciones de recorte en
\cref{eq:angulos_saturacion_sesgada} hacen que
\cref{eq:psi0_saturacion_cerrada,eq:psi1_saturacion_cerrada} sigan siendo
válidas cuando alguna de las tres regiones tiene longitud cero.

La ecuación de media puede reducirse a una igualdad escalar. De
\cref{eq:transferencia_cramer},
\begin{align}
W_q(0)
&=-\alpha
\frac{\beta+\gamma}
{\alpha(1+\ell)(\beta+\gamma)-\alpha\gamma}\nonumber\\
&=-\frac{\beta+\gamma}
{(1+\ell)(\beta+\gamma)-\gamma}.
\label{eq:ganancia_dc}
\end{align}
Por tanto,
\begin{equation}
c=W_q(0)\Psi_0(c,A),
\label{eq:cierre_dc_escalar}
\end{equation}
y el cierre armónico es
\[
1-W_q(\ii\omega)\frac{\Psi_1(c,A)}{A}=0.
\]
Estas tres ecuaciones reales ---una para la media, y las partes real e
imaginaria del cierre--- determinan \((c,A,\omega)\).

\fi
\ifnum\HAFOderivationpart=8\relax
\subsection{Derivación completa para la arctangente}
\label{app:df_arctan}

Sea
\[
\psi_{\mathrm a}(\sigma)=a_2\arctan(\rho\sigma),
\qquad
\xi=\rho A.
\]
Se define
\[
I(\xi)=
\int_0^\pi\arctan(\xi\sin\theta)\sin\theta\dd\theta.
\]
Como \(I(0)=0\), puede calcularse \(I\) integrando su derivada:
\begin{align}
I'(\xi)
&=\int_0^\pi
\frac{\sin^2\theta}{1+\xi^2\sin^2\theta}\dd\theta\nonumber\\
&=\frac1{\xi^2}
\left[
\pi-\int_0^\pi
\frac{\dd\theta}{1+\xi^2\sin^2\theta}
\right]\nonumber\\
&=\frac{\pi}{\xi^2}
\left(1-\frac1{\sqrt{1+\xi^2}}\right).
\label{eq:derivada_integral_arctan}
\end{align}
En la última igualdad se usó la integral estándar
\[
\int_0^\pi\frac{\dd\theta}{1+\xi^2\sin^2\theta}
=\frac{\pi}{\sqrt{1+\xi^2}}.
\]
Una primitiva que se anula en \(\xi=0\) es
\begin{equation}
I(\xi)
=\pi\frac{\sqrt{1+\xi^2}-1}{\xi}.
\label{eq:integral_arctan}
\end{equation}
Al sustituir \(\xi=\rho A\) en
\cref{eq:df_centrada},
\begin{align}
N_{\mathrm a}(A)
&=\frac{2a_2}{\pi A}I(\rho A)\nonumber\\
&=\frac{2a_2}{\rho A^2}
\left(\sqrt{1+\rho^2A^2}-1\right)\nonumber\\
&=\frac{2a_2\rho}
{1+\sqrt{1+\rho^2A^2}}.
\label{eq:df_arctan_derivada}
\end{align}

Si el cierre exige \(N_{\mathrm a}(A)=k\), la amplitud puede despejarse.
De la forma racionalizada,
\[
1+\sqrt{1+\rho^2A^2}=\frac{2a_2\rho}{k}.
\]
Al aislar la raíz, elevar al cuadrado y simplificar,
\begin{equation}
A^2
=\frac{4a_2(a_2\rho-k)}{k^2\rho}.
\label{eq:amplitud_arctan}
\end{equation}
La solución se acepta solamente si el lado derecho es positivo y la igualdad
previa a elevar al cuadrado tiene lado derecho no menor que \(2\). De manera
equivalente, para \(A>0\), \(k\) tiene el mismo signo que \(a_2\) y se
encuentra estrictamente entre \(0\) y \(a_2\rho\), con el orden de los
extremos ajustado al signo de \(a_2\).

\fi
\ifnum\HAFOderivationpart=9\relax
\subsection{Derivación completa de la identidad modal y su normalización}
\label{app:identidad_modal}

Sea \(P_0=P+bkr^{\T}\). Entonces
\[
\zeta I-P_0=(\zeta I-P)-bkr^{\T}.
\]
Si \(\zeta I-P\) es invertible, el lema del determinante matricial
\(\det(M+uv^{\T})=\det(M)(1+v^{\T}M^{-1}u)\), con
\[
M=\zeta I-P,\qquad u=-bk,\qquad v=r,
\]
produce
\begin{align}
\det(\zeta I-P_0)
&=\det(\zeta I-P)
\left[1-kr^{\T}(\zeta I-P)^{-1}b\right]\nonumber\\
&=\det(\zeta I-P)\bigl[1-kG(\zeta)\bigr],
\label{eq:lema_determinante_aplicado}
\end{align}
donde
\[
G(\zeta)=r^{\T}(\zeta I-P)^{-1}b,
\qquad W_q(s)=G(s^q).
\]
Por continuidad, la identidad polinómica se extiende también a los valores en
los que la resolvente de \(P\) es singular. Si
\[
1-kG(\zeta_0)=0,
\]
entonces \(\det(\zeta_0I-P_0)=0\), de modo que existe
\(v\neq0\) con \(P_0v=\zeta_0v\). Siempre que \(r^{\T}v\neq0\), se reemplaza
\(v\) por \(v/(r^{\T}v)\) y se obtiene \(r^{\T}v=1\). Ésta es la
normalización usada en \cref{eq:modo_normalizado,eq:semilla_modal}.

En la rama sesgada,
\[
X_1=\bigl[(\ii\omega_0)^qI-P\bigr]^{-1}b\Psi_1.
\]
Al premultiplicar por \(r^{\T}\),
\[
r^{\T}X_1
=W_q(\ii\omega_0)\Psi_1
=W_q(\ii\omega_0)N_1A.
\]
El cierre \(1-W_qN_1=0\) prueba algebraicamente que
\(r^{\T}X_1=A\).

\fi
\ifnum\HAFOderivationpart=10\relax
\subsection{Derivación completa de la homotopía afín}
\label{app:homotopia_afin}

La aproximación sesgada separa
\[
\psi(\sigma)
=\Psi_0+N_1(\sigma-c)
+\bigl[\psi(\sigma)-\Psi_0-N_1(\sigma-c)\bigr].
\]
Se definen
\begin{equation}
P_{\mathrm a}=P+N_1br^{\T},\qquad
d_{\mathrm a}=b(\Psi_0-N_1c).
\label{eq:parametros_homotopia_afin}
\end{equation}
La familia de campos es
\begin{equation}
\begin{aligned}
F_\varepsilon(X)
&=P_{\mathrm a}X+d_{\mathrm a}\\
&\quad+\varepsilon b
\left[
\psi(r^{\T}X)-\Psi_0-N_1(r^{\T}X-c)
\right],
\qquad 0\leq\varepsilon\leq1.
\end{aligned}
\label{eq:homotopia_afin}
\end{equation}

En \(\varepsilon=0\), la media \(\bar X\) de
\cref{eq:cierre_dc_sesgado} es un equilibrio del auxiliar:
\begin{align*}
F_0(\bar X)
&=P\bar X+bN_1c+b\Psi_0-bN_1c\\
&=P\bar X+b\Psi_0=0.
\end{align*}
Para una perturbación \(\xi=X-\bar X\),
\[
F_0(\bar X+\xi)=P_{\mathrm a}\xi,
\]
de modo que el auxiliar conserva tanto la media como el modo armónico.

En \(\varepsilon=1\), al agrupar términos,
\begin{align*}
F_1(X)
&=PX+bN_1r^{\T}X+b\Psi_0-bN_1c\\
&\quad+b\psi(r^{\T}X)-b\Psi_0-bN_1r^{\T}X+bN_1c\\
&=PX+b\psi(r^{\T}X).
\end{align*}
Así, el extremo final de \cref{eq:homotopia_afin} es exactamente el sistema
físico de \cref{eq:lure_comun}. Una homotopía centrada de la forma
\(P_0X+\varepsilon b[\psi(\sigma)-k\sigma]\) no conserva, en general, la
media \(c\neq0\); por eso no se utiliza para la rama sesgada.
\fi

\endgroup
Como
\(W_q(\ii\omega_0)=G((\ii\omega_0)^q)=G(\zeta_0)\), al cumplirse el cierre
\(\zeta_0\) es un valor propio del sistema lineal auxiliar. Se calcula un
vector \(v\) mediante
\begin{equation}
(P_0-\zeta_0I)v=0,\qquad r^{\T}v=1,
\label{eq:modo_normalizado}
\end{equation}
y la familia de fases centradas es
\begin{equation}
X_{\mathrm{seed}}(\theta_s)
=A_0\Re\!\left(\e^{\ii\theta_s}v\right),
\qquad 0\leq\theta_s<2\pi.
\label{eq:semilla_modal}
\end{equation}
Además del residuo de \cref{eq:modo_normalizado}, se verifica que el tercer
modo de \(P_0\) queda en el sector estable de Matignon. El par
\(\zeta_0,\overline{\zeta_0}\) se encuentra sobre la frontera angular
\(|\arg\zeta|=q\pi/2\) del problema auxiliar; ésa es la contraparte
fraccionaria del par \(\pm\ii\omega_0\) del caso entero.

Para la rama sesgada se reconstruyen por separado la media y el armónico. Con
\(\bar X\) dada por \cref{eq:cierre_dc_sesgado},
\begin{equation}
\begin{aligned}
\bigl[(\ii\omega_0)^qI-P\bigr]X_1
&=b\,\Psi_1(c,A),\\
X_{\mathrm{seed}}(\theta_s)
&=\bar X+\Re\!\left(\e^{\ii\theta_s}X_1\right).
\end{aligned}
\label{eq:semilla_sesgada}
\end{equation}
El cierre implica
\[
r^{\T}X_1=W_q(\ii\omega_0)\Psi_1
=W_q(\ii\omega_0)N_1A=A,
\]
de modo que la primera componente de
\cref{eq:semilla_sesgada} tiene exactamente la media \(c\) y la amplitud
\(A\) usadas en \cref{eq:df_sesgada}.

La semilla armónica pertenece al problema auxiliar de Liouville--Weyl. La
trayectoria física se obtiene después al integrarla y continuarla con el
operador causal de Caputo. Por tanto, el cierre armónico es un procedimiento
de localización, no una sustitución de la integración causal.

\subsection{Defecto de arranque entre el armónico de Weyl y el PVI de Caputo}
\label{subsec:armonic-weyl-caputo-defect}

La identidad estacionaria
\(D_{\mathrm W}^{q}e^{\ii\omega t}=(\ii\omega)^qe^{\ii\omega t}\)
utiliza una historia extendida hacia \(-\infty\). Para \(0<q<1\) y
\(\omega>0\), puede justificarse por regularización de Abel:
\begin{equation}
 \int_0^\infty u^{-q}e^{-(\varepsilon+\ii\omega)u}\,\dd u
 =\Gamma(1-q)(\varepsilon+\ii\omega)^{q-1},
 \qquad\varepsilon>0.
 \label{eq:weyl-abel-harmonic}
\end{equation}
Al multiplicar por \(\ii\omega/\Gamma(1-q)\) y tomar
\(\varepsilon\downarrow0\), se obtiene el multiplicador
\((\ii\omega)^q\) en la rama principal. La formulación por Fourier
produce el mismo resultado.

En cambio, para el armónico suave iniciado en \(t_0=0\), la definición
de Caputo y el cambio \(u=t-s\) dan
\begin{equation}
 {}^{\mathrm C}D_0^q e^{\ii\omega t}
 =\frac{\ii\omega e^{\ii\omega t}}{\Gamma(1-q)}
    \int_0^t u^{-q}e^{-\ii\omega u}\,\dd u.
 \label{eq:caputo-harmonic-finite-history}
\end{equation}
Restar la integral estacionaria muestra explícitamente el defecto:
\begin{equation}
 \begin{aligned}
 {}^{\mathrm C}D_0^q e^{\ii\omega t}
       -(\ii\omega)^qe^{\ii\omega t}
 &=-\frac{\ii\omega e^{\ii\omega t}}{\Gamma(1-q)}\\
 &\quad\times\int_t^\infty u^{-q}e^{-\ii\omega u}\,\dd u.
 \end{aligned}
 \label{eq:caputo-weyl-startup-defect}
\end{equation}
La integral impropia converge por el criterio de Dirichlet. El lado de
Caputo satisface, cerca de cero,
\[
 {}^{\mathrm C}D_0^q e^{\ii\omega t}
 =\frac{\ii\omega}{\Gamma(2-q)}t^{1-q}
      +O(t^{2-q}),
\]
y tiende a cero, mientras que
\((\ii\omega)^qe^{\ii\omega t}\to(\ii\omega)^q\ne0\).
Esta comparación corresponde a un armónico suave prescrito, no a la
regularidad de toda solución de un PVI fraccionario. Una solución causal
del auxiliar requiere la respuesta de arranque impuesta por su dato inicial.

El truncamiento armónico introduce un defecto diferente: las componentes
de Fourier de \(\psi\) que no se retienen en el cierre. Ni la igualdad
de la media ni la del fundamental anulan esas componentes. Por eso, la
semilla se obtiene del problema estacionario y se evalúa después en el
PVI causal, sin identificar ambos problemas. Los resultados de
no periodicidad con terminal inferior finito respaldan esta separación
\cite{TavazoeiHaeri2009,AreaLosadaNieto2014}.

Danca describe la linealización armónica y la continuación como métodos
de localización, y resuelve los PVI fraccionarios mediante ABM
\cite{Danca2017}. Wu et al. trasladan el algoritmo de Leonov a Chua
con arctangente y construyen un auxiliar con transferencias racionales
en una variable compleja \(p\), incluido un bloque
\(p^2+\omega_0^2\) \cite{Wu2023ArctanChua}. Una similitud matricial
sigue siendo válida para Caputo, pero interpretar los polos como una
frecuencia temporal requiere la relación espectral
\(p=(\ii\omega)^q\). Para \(q<1\), un par conservado en
\(\pm\ii\omega_0\) queda en el sector estable, mientras que
\cref{eq:matriz_auxiliar} sitúa el par en la frontera angular.

Wu et al. reconocen expresamente la no periodicidad exacta y la pérdida
de parte de la memoria histórica en su aproximación por descomposición
de Adomian en subintervalos; también distinguen los contornos
numéricamente periódicos de la periodicidad matemática
\cite{Wu2023ArctanChua}. Reutilizar una semilla entera es un ensayo
causal legítimo, pero no demuestra la persistencia del objeto que la
motivó. La función descriptiva estática puede conservarse sin introducir
\(q\) en su integral de Fourier; la modificación del bloque dinámico
y el tratamiento de la historia deben justificarse por separado.

\subsection{Recuperación del caso entero y separación del caso \texttt{c590}}

Para \(X\in C^1\), con \(t>t_0\), y valores donde existe la resolvente
en la rama fijada, se tienen los límites siguientes; el del núcleo se
entiende para \(\tau<t\):
\begin{equation}
\begin{aligned}
\lim_{q\to1^-}{}^{\mathrm C}D_{t_0}^{q}X(t)&=\dot X(t),&
\lim_{q\to1^-}(\ii\omega)^q&=\ii\omega,\\
\lim_{q\to1^-}W_q(s)
&=r^{\T}(sI-P)^{-1}b,&
\lim_{q\to1^-}\frac{(t-\tau)^{q-1}}{\Gamma(q)}&=1.
\end{aligned}
\label{eq:limite_q_uno}
\end{equation}
Las matrices de la realización y las integrales estáticas de Fourier no
dependen del orden. La continuidad del inverso matricial, que se deduce
de su expresión mediante adjunta y determinante no nulo, hace que
\cref{eq:cierre_armonico,eq:modo_normalizado,eq:semilla_modal} recuperen
las fórmulas enteras cuando \(q=1\). La continuidad de las raíces del
balance requiere, además, una condición de no degeneración.

Para la rama sesgada, defina el residual real
\begin{equation}
 \mathcal F(c,A,\omega;q)=
 \begin{pmatrix}
 c-W_q(0)\Psi_0(c,A)\\
 \Re[1-W_q(\ii\omega)N_1(c,A)]\\
 \Im[1-W_q(\ii\omega)N_1(c,A)]
 \end{pmatrix}.
 \label{eq:balance-branch-implicit-residual}
\end{equation}
Si \(\mathcal F\) es \(C^1\) cerca de una raíz con \(q=1\),
\(A>0\), \(\omega>0\), y
\(D_{(c,A,\omega)}\mathcal F\) es invertible, el teorema de la función
implícita proporciona una rama local
\((c(q),A(q),\omega(q))\) que converge a esa raíz. Para la rama
centrada se usa el residual de dos componentes reales en \((A,\omega)\).
Se mantienen fuera de los polos la frecuencia y la ganancia continua
cuando ésta interviene.

En dicha rama, la ganancia, la media y el modo normalizado son funciones
continuas de \(q\); por tanto, también lo son la semilla y los
coeficientes de la homotopía afín. En \(q=1\),
\(\bar X+\Re(Ave^{\ii\omega t})\) es una solución exacta del
auxiliar, pues \(P_0v=\ii\omega v\) y la media es un equilibrio
del campo afín. El campo no lineal final conserva los errores de
truncamiento hasta que se estudia su trayectoria. Si el Jacobiano del
balance se hace singular, la rama puede bifurcarse o terminar; los
límites de las fórmulas no prueban continuación global ni persistencia
de un atractor \cite{LeonovKuznetsov2013,KuznetsovEtAl2017}.

La fórmula \cref{eq:df_arctan} corresponde a la no linealidad bibliográfica con
\(\rho=1\) y también la forma algebraica general con \(\rho>0\). Sin embargo,
la condición inicial del caso \texttt{c590} de
\cref{eq:parametros_c590} procede de una búsqueda acotada en orden entero y
de su refinamiento mediante problemas de Caputo independientes. Esta
construcción no utiliza el cierre armónico de arctangente.
